% Options for packages loaded elsewhere
\PassOptionsToPackage{unicode}{hyperref}
\PassOptionsToPackage{hyphens}{url}
%
\documentclass[
]{article}
\usepackage{amsmath,amssymb}
\usepackage{iftex}
\ifPDFTeX
  \usepackage[T1]{fontenc}
  \usepackage[utf8]{inputenc}
  \usepackage{textcomp} % provide euro and other symbols
\else % if luatex or xetex
  \usepackage{unicode-math} % this also loads fontspec
  \defaultfontfeatures{Scale=MatchLowercase}
  \defaultfontfeatures[\rmfamily]{Ligatures=TeX,Scale=1}
\fi
\usepackage{lmodern}
\ifPDFTeX\else
  % xetex/luatex font selection
\fi
% Use upquote if available, for straight quotes in verbatim environments
\IfFileExists{upquote.sty}{\usepackage{upquote}}{}
\IfFileExists{microtype.sty}{% use microtype if available
  \usepackage[]{microtype}
  \UseMicrotypeSet[protrusion]{basicmath} % disable protrusion for tt fonts
}{}
\makeatletter
\@ifundefined{KOMAClassName}{% if non-KOMA class
  \IfFileExists{parskip.sty}{%
    \usepackage{parskip}
  }{% else
    \setlength{\parindent}{0pt}
    \setlength{\parskip}{6pt plus 2pt minus 1pt}}
}{% if KOMA class
  \KOMAoptions{parskip=half}}
\makeatother
\usepackage{xcolor}
\usepackage{color}
\usepackage{fancyvrb}
\newcommand{\VerbBar}{|}
\newcommand{\VERB}{\Verb[commandchars=\\\{\}]}
\DefineVerbatimEnvironment{Highlighting}{Verbatim}{commandchars=\\\{\}}
% Add ',fontsize=\small' for more characters per line
\newenvironment{Shaded}{}{}
\newcommand{\AlertTok}[1]{\textcolor[rgb]{1.00,0.00,0.00}{\textbf{#1}}}
\newcommand{\AnnotationTok}[1]{\textcolor[rgb]{0.38,0.63,0.69}{\textbf{\textit{#1}}}}
\newcommand{\AttributeTok}[1]{\textcolor[rgb]{0.49,0.56,0.16}{#1}}
\newcommand{\BaseNTok}[1]{\textcolor[rgb]{0.25,0.63,0.44}{#1}}
\newcommand{\BuiltInTok}[1]{\textcolor[rgb]{0.00,0.50,0.00}{#1}}
\newcommand{\CharTok}[1]{\textcolor[rgb]{0.25,0.44,0.63}{#1}}
\newcommand{\CommentTok}[1]{\textcolor[rgb]{0.38,0.63,0.69}{\textit{#1}}}
\newcommand{\CommentVarTok}[1]{\textcolor[rgb]{0.38,0.63,0.69}{\textbf{\textit{#1}}}}
\newcommand{\ConstantTok}[1]{\textcolor[rgb]{0.53,0.00,0.00}{#1}}
\newcommand{\ControlFlowTok}[1]{\textcolor[rgb]{0.00,0.44,0.13}{\textbf{#1}}}
\newcommand{\DataTypeTok}[1]{\textcolor[rgb]{0.56,0.13,0.00}{#1}}
\newcommand{\DecValTok}[1]{\textcolor[rgb]{0.25,0.63,0.44}{#1}}
\newcommand{\DocumentationTok}[1]{\textcolor[rgb]{0.73,0.13,0.13}{\textit{#1}}}
\newcommand{\ErrorTok}[1]{\textcolor[rgb]{1.00,0.00,0.00}{\textbf{#1}}}
\newcommand{\ExtensionTok}[1]{#1}
\newcommand{\FloatTok}[1]{\textcolor[rgb]{0.25,0.63,0.44}{#1}}
\newcommand{\FunctionTok}[1]{\textcolor[rgb]{0.02,0.16,0.49}{#1}}
\newcommand{\ImportTok}[1]{\textcolor[rgb]{0.00,0.50,0.00}{\textbf{#1}}}
\newcommand{\InformationTok}[1]{\textcolor[rgb]{0.38,0.63,0.69}{\textbf{\textit{#1}}}}
\newcommand{\KeywordTok}[1]{\textcolor[rgb]{0.00,0.44,0.13}{\textbf{#1}}}
\newcommand{\NormalTok}[1]{#1}
\newcommand{\OperatorTok}[1]{\textcolor[rgb]{0.40,0.40,0.40}{#1}}
\newcommand{\OtherTok}[1]{\textcolor[rgb]{0.00,0.44,0.13}{#1}}
\newcommand{\PreprocessorTok}[1]{\textcolor[rgb]{0.74,0.48,0.00}{#1}}
\newcommand{\RegionMarkerTok}[1]{#1}
\newcommand{\SpecialCharTok}[1]{\textcolor[rgb]{0.25,0.44,0.63}{#1}}
\newcommand{\SpecialStringTok}[1]{\textcolor[rgb]{0.73,0.40,0.53}{#1}}
\newcommand{\StringTok}[1]{\textcolor[rgb]{0.25,0.44,0.63}{#1}}
\newcommand{\VariableTok}[1]{\textcolor[rgb]{0.10,0.09,0.49}{#1}}
\newcommand{\VerbatimStringTok}[1]{\textcolor[rgb]{0.25,0.44,0.63}{#1}}
\newcommand{\WarningTok}[1]{\textcolor[rgb]{0.38,0.63,0.69}{\textbf{\textit{#1}}}}
\usepackage{longtable,booktabs,array}
\usepackage{calc} % for calculating minipage widths
% Correct order of tables after \paragraph or \subparagraph
\usepackage{etoolbox}
\makeatletter
\patchcmd\longtable{\par}{\if@noskipsec\mbox{}\fi\par}{}{}
\makeatother
% Allow footnotes in longtable head/foot
\IfFileExists{footnotehyper.sty}{\usepackage{footnotehyper}}{\usepackage{footnote}}
\makesavenoteenv{longtable}
\setlength{\emergencystretch}{3em} % prevent overfull lines
\providecommand{\tightlist}{%
  \setlength{\itemsep}{0pt}\setlength{\parskip}{0pt}}
\setcounter{secnumdepth}{-\maxdimen} % remove section numbering
\ifLuaTeX
  \usepackage{selnolig}  % disable illegal ligatures
\fi
\usepackage{bookmark}
\IfFileExists{xurl.sty}{\usepackage{xurl}}{} % add URL line breaks if available
\urlstyle{same}
\hypersetup{
  pdftitle={Ratcheting Consequence Inquiry},
  hidelinks,
  pdfcreator={LaTeX via pandoc}}

\title{Ratcheting Consequence Inquiry}
\usepackage{etoolbox}
\makeatletter
\providecommand{\subtitle}[1]{% add subtitle to \maketitle
  \apptocmd{\@title}{\par {\large #1 \par}}{}{}
}
\makeatother
\subtitle{A Question-Driven, Proof-Producing Architecture for
Domain-Agnostic Inquiry}
\author{}
\date{}

\begin{document}
\maketitle

\section{Ratcheting Consequence
Inquiry}\label{ratcheting-consequence-inquiry}

\subsection{A Question-Driven, Proof-Producing Architecture for
Domain-Agnostic
Inquiry}\label{a-question-driven-proof-producing-architecture-for-domain-agnostic-inquiry}

\textbf{Suggested filename:} \texttt{RCI\_Project\_Spec.md}\\
\textbf{Status:} Project specification / formal design document\\
\textbf{Version:} 0.1

\begin{center}\rule{0.5\linewidth}{0.5pt}\end{center}

\subsection{Abstract}\label{abstract}

Ratcheting Consequence Inquiry (RCI) is a domain-agnostic inquiry
architecture in which questions are the lawful control language of
reasoning.

RCI separates four things that are usually conflated:

\begin{enumerate}
\def\labelenumi{\arabic{enumi}.}
\tightlist
\item
  semantic generation --- proposing distinctions, examples,
  transformations, explanations, counterexamples, prerequisites, and
  other candidate content;
\item
  question structure --- determining what kind of relation a semantic
  answer is being asked to instantiate;
\item
  formal control --- composing claims, testing consequences, refining
  abstractions, generating obligations, localizing conflicts, and
  selecting subsequent questions;
\item
  warrant --- deciding which provisional claims may be promoted into
  retained knowledge.
\end{enumerate}

The central architectural move is:

\[
\boxed{
Q_\tau
:
\Sigma
\rightsquigarrow
\operatorname{Claim}[\tau]
}
\]

rather than

\[
Q_\tau:\Sigma\rightsquigarrow\tau.
\]

A question determines the type of claim requested. The answer supplies
an arbitrary semantic payload. The controller need not initially
understand or verify that payload deeply. It provisionally treats it
according to the relation represented by the question that elicited it.

Thus:

\[
\boxed{
\text{question}
\to
\text{typed claim role}
\to
\text{semantic payload}
\to
\text{lawful composition}
\to
\text{conflict / support / boundary / obligation}
\to
\text{next question}.
}
\]

An incorrect semantic answer does not make the inquiry protocol
ill-typed. It creates a possibly false claim. Subsequent lawful
questions are designed to expose false necessities, false sufficiencies,
hidden guards, inconsistent commitments, arrangement conflicts,
succession failures, inadequate abstractions, and unsupported
generalizations.

The architecture therefore does not require the language model to be
reliably correct at every step. It requires:

\[
\boxed{
\textbf{fallibility to be composable.}
}
\]

RCI combines a general relational theory with established machinery from
abstract interpretation, predicate abstraction, CEGAR, active automata
learning, IC3/PDR, DPLL(T), constrained Horn clauses, predicate
transformers, conflict minimization, assumption-based truth maintenance,
interpolation, and controller refinement. Where no single established
formalism spans the whole system, this document defines the intermediate
contracts required to compose them lawfully.

The result is simultaneously:

\begin{itemize}
\tightlist
\item
  a general theory of inquiry;
\item
  an executable reasoning architecture;
\item
  a typed question calculus;
\item
  a protocol for using untrusted semantic generators;
\item
  and an implementation specification.
\end{itemize}

\begin{center}\rule{0.5\linewidth}{0.5pt}\end{center}

\section{Part I --- Project Thesis}\label{part-i-project-thesis}

\subsection{1. Core claim}\label{core-claim}

Inquiry advances by discovering which distinctions change protected
consequences and eliminating distinctions that do not.

The two primitive epistemic movements are:

\[
\boxed{
\text{remove unsupported possibilities}
}
\]

and

\[
\boxed{
\text{remove false equivalences}.
}
\]

The central ratchet is:

\[
\boxed{
\textbf{Never retain a reasoning step merely because content was generated.
Retain it only if it produces a warranted consequential change.}
}
\]

Candidate generation may expand arbitrarily.

Retained inquiry must ratchet.

\begin{center}\rule{0.5\linewidth}{0.5pt}\end{center}

\subsection{2. Questions as control
operators}\label{questions-as-control-operators}

Natural-language questions are treated as executable control operators
over an evolving inquiry state.

A question does not merely request information.

It establishes:

\begin{enumerate}
\def\labelenumi{\arabic{enumi}.}
\tightlist
\item
  which prior objects are being related;
\item
  which relation is being requested;
\item
  which kind of answer may lawfully be created;
\item
  which subsequent operations may consume that answer.
\end{enumerate}

For example, the surface question

\begin{quote}
Can the result still hold without this?
\end{quote}

corresponds formally to a necessity-counterexample obligation.

The controller already knows:

\begin{itemize}
\tightlist
\item
  the protected consequence \(C\);
\item
  the proposed necessity \(N\);
\item
  that a positive answer is being offered as a candidate member of
  \(C\land\neg N\).
\end{itemize}

The language model need not express the logical form itself.

Thus:

\[
\boxed{
\textbf{the question carries structural semantics;
the answer supplies local semantic content.}
}
\]

\begin{center}\rule{0.5\linewidth}{0.5pt}\end{center}

\subsection{3. Trusted control, untrusted
payload}\label{trusted-control-untrusted-payload}

The central separation is:

\[
\boxed{
\text{lawful question composition}
\neq
\text{truth of answer payload}.
}
\]

A generated answer may be:

\begin{itemize}
\tightlist
\item
  correct;
\item
  incomplete;
\item
  vague;
\item
  overly broad;
\item
  contradictory;
\item
  mis-scoped;
\item
  semantically incoherent;
\item
  or wrong.
\end{itemize}

The controller nevertheless remains well-defined because the answer
initially becomes:

\[
\operatorname{Claim}[\tau],
\]

not an established object of type \(\tau\).

Only later may it become:

\[
\operatorname{Candidate}[\tau],
\]

and eventually:

\[
\operatorname{Warranted}[\tau].
\]

\begin{center}\rule{0.5\linewidth}{0.5pt}\end{center}

\subsection{4. The project objective}\label{the-project-objective}

The project is to build a system in which a general semantic model
supplies candidate content while a formal controller:

\begin{itemize}
\tightlist
\item
  determines what relational object should be sought next;
\item
  compiles that need into a plain-language question;
\item
  records the answer under the type implied by that question;
\item
  composes the resulting claim with earlier commitments;
\item
  deliberately seeks counterexamples and contradictions;
\item
  localizes whatever fails;
\item
  asks the next question implied by that failure;
\item
  invokes formal or external machinery only when needed;
\item
  generalizes warranted returns;
\item
  compresses irrelevant distinctions;
\item
  and recursively repeats the process.
\end{itemize}

The intended outer loop is:

\[
\boxed{
\text{obligation}
\to
\text{question}
\to
\text{claim}
\to
\text{composition}
\to
\text{return}
\to
\text{warrant}
\to
\text{update}
\to
\text{next obligation}.
}
\]

\begin{center}\rule{0.5\linewidth}{0.5pt}\end{center}

\section{Part II --- Formal Substrate}\label{part-ii-formal-substrate}

\subsection{5. Typed binding}\label{typed-binding}

An inquiry is conducted relative to a binding

\[
\boxed{
\mathfrak B
=
(\mathbf X,\mathcal L,\mathcal A,\mathcal U,\mathcal H,\mathcal E)
}
\]

where:

\begin{itemize}
\tightlist
\item
  \(\mathbf X=\{X_\tau\}\) is the family of typed carriers;
\item
  \(\mathcal L\) is the available relational language;
\item
  \(\mathcal A\) is the family of admissible arrangement deformations;
\item
  \(\mathcal U\) is the family of admissible successions;
\item
  \(\mathcal H\) is the protected consequence horizon;
\item
  \(\mathcal E\) is the family of available independent discharge
  mechanisms.
\end{itemize}

An inquiry does not assume the existence of:

\begin{itemize}
\tightlist
\item
  a metric;
\item
  probability;
\item
  temporal order;
\item
  determinism;
\item
  causality;
\item
  scalar utility;
\item
  continuity;
\item
  total comparability;
\item
  or complete observability.
\end{itemize}

Such structure may be added only when warranted by the binding.

\begin{center}\rule{0.5\linewidth}{0.5pt}\end{center}

\subsection{6. Typed relations}\label{typed-relations}

For carriers \(X,Y\in\mathbf X\), a relation is

\[
R:X\rightsquigarrow Y
\]

with

\[
R\subseteq X\times Y.
\]

Relations compose:

\[
S\circ R
=
\{(x,z):\exists y,\ xRy\land ySz\}.
\]

For a relation

\[
R:X\rightsquigarrow Y
\]

and a target predicate \(Q\subseteq Y\), define the existential
predecessor:

\[
\boxed{
\Diamond_R(Q)
=
\{x:R[x]\cap Q\neq\varnothing\}
}
\]

and universal predecessor:

\[
\boxed{
\Box_R(Q)
=
\{x:R[x]\subseteq Q\}.
}
\]

When execution must exist:

\[
\boxed{
\operatorname{Must}_R(Q)
=
\operatorname{Dom}(R)\cap\Box_R(Q).
}
\]

These satisfy:

\[
\Diamond_{S\circ R}
=
\Diamond_R\circ\Diamond_S
\]

and

\[
\Box_{S\circ R}
=
\Box_R\circ\Box_S.
\]

This is the mathematical basis of succession, backward prerequisite
propagation, reachability, and actualization.

Kleene algebra with tests supplies an established algebraic basis for
compositional reasoning over program-like actions and embedded
predicates. {[}R1{]}

\begin{center}\rule{0.5\linewidth}{0.5pt}\end{center}

\section{Part III --- Protected
Consequence}\label{part-iii-protected-consequence}

\subsection{7. Consequence profile}\label{consequence-profile}

Let

\[
\Phi:X\rightsquigarrow Z
\]

be the current protected consequence profile.

A focal binary consequence is

\[
C:X\to\{0,1\}.
\]

Define

\[
C^+=C^{-1}(1),
\qquad
C^-=C^{-1}(0).
\]

The consequence polarity is relational rather than intrinsically
normative. Either side may be desired, undesired, neutral, unknown, or
merely contrasting.

\begin{center}\rule{0.5\linewidth}{0.5pt}\end{center}

\subsection{8. Problem and non-problem}\label{problem-and-non-problem}

If the inquiry begins from a recognized problem predicate

\[
P\subseteq X,
\]

then the first opposite class is simply

\[
\boxed{
\neg P.
}
\]

No concrete solution exemplar is required.

However:

\[
\neg P
\]

is only the provisional non-problem class.

Further protected consequences may distinguish acceptable resolution
from merely absence of the original problem.

Thus:

\[
\boxed{
\text{problem recognition supplies an opposite pole,
not necessarily a complete solution specification.}
}
\]

\begin{center}\rule{0.5\linewidth}{0.5pt}\end{center}

\section{Part IV --- Consequence Equivalence and
Distinction}\label{part-iv-consequence-equivalence-and-distinction}

\subsection{9. Consequence equivalence}\label{consequence-equivalence}

For protected continuations \(h\in\mathcal H\), let

\[
J_h:X\rightsquigarrow Z_h.
\]

Define:

\[
\boxed{
x\equiv_{\mathcal H}y
\iff
\forall h\in\mathcal H,
\quad
J_h[x]=J_h[y].
}
\]

Therefore:

\[
x\neq y
\]

does not imply

\[
x\not\equiv_{\mathcal H}y.
\]

The consequence-relative state space is

\[
\boxed{
X/{\equiv_{\mathcal H}}.
}
\]

This is the formal basis of RCI compression.

For deterministic sequential systems, this resembles the
continuation-based equivalence underlying Myhill--Nerode; Angluin's
\(L^*\) algorithm actively learns such behaviorally distinct states by
membership/equivalence queries and counterexamples. {[}R2{]}

For explicit relational systems, efficient partition-refinement methods
such as Paige--Tarjan compute coarsest relational partitions. {[}R3{]}

\begin{center}\rule{0.5\linewidth}{0.5pt}\end{center}

\subsection{10. Elementary distinction}\label{elementary-distinction}

Let

\[
d:X\to D.
\]

If

\[
d(x)\neq d(y),
\]

then the elementary distinction is:

\[
\boxed{
(x,d,y).
}
\]

This is triadic:

\begin{itemize}
\tightlist
\item
  first relatum;
\item
  second relatum;
\item
  common relation under which they differ.
\end{itemize}

The triad is not an additional algorithmic primitive. It is the minimum
geometry of a distinction.

\begin{center}\rule{0.5\linewidth}{0.5pt}\end{center}

\subsection{11. Tetradic refinement}\label{tetradic-refinement}

Introduce a second binary distinction

\[
D:X\to\{0,1\}.
\]

The joint refinement is:

\[
\boxed{
\langle C,D\rangle
:
X\to\{0,1\}^2.
}
\]

Its realized image is:

\[
\boxed{
T_{C,D}
=
\operatorname{im}\langle C,D\rangle.
}
\]

The logical cells are:

\[
C\land D,
\]

\[
C\land\neg D,
\]

\[
\neg C\land D,
\]

\[
\neg C\land\neg D.
\]

The existence of two distinctions does not imply that all four cells are
realizable.

A missing cell is an inquiry target.

The recursive geometry is therefore:

\[
\boxed{
\text{triad}
\to
\text{second distinction}
\to
\text{tetradic refinement}
\to
\text{residual}
\to
\text{new triad}.
}
\]

\begin{center}\rule{0.5\linewidth}{0.5pt}\end{center}

\section{Part V --- Arrangement and
Succession}\label{part-v-arrangement-and-succession}

\subsection{12. Arrangement}\label{arrangement}

Arrangement concerns differences among realizations.

For a reference arrangement \(r\), let

\[
\preceq_A^r
\]

be a binding-supplied preorder.

\[
x\preceq_A^r y
\]

means that \(x\) preserves at least as much relevant arrangement
structure of \(r\) as \(y\).

No numerical distance is required.

\begin{center}\rule{0.5\linewidth}{0.5pt}\end{center}

\subsection{13. Succession}\label{succession}

Each

\[
u\in\mathcal U
\]

is represented as

\[
R_u\subseteq X\times X.
\]

A succession path

\[
\pi=(u_1,\ldots,u_n)
\]

has relation:

\[
R_\pi
=
R_{u_n}\circ\cdots\circ R_{u_1}.
\]

A binding may supply a path preorder:

\[
\preceq_S.
\]

Arrangement minimality and succession minimality are distinct.

\[
\boxed{
\operatorname{MinArrangement}
\neq
\operatorname{MinSuccession}.
}
\]

\begin{center}\rule{0.5\linewidth}{0.5pt}\end{center}

\section{Part VI --- The Central Inquiry
Geometry}\label{part-vi-the-central-inquiry-geometry}

\subsection{14. Stretch within a consequence
class}\label{stretch-within-a-consequence-class}

The first geometric pressure is:

\[
\boxed{
\text{maximize variation while consequence remains fixed.}
}
\]

Given current predicate basis \(B\),

\[
\Delta_B(x,y)
=
\{p\in B:p(x)\neq p(y)\}.
\]

For \(x\in C^b\),

\[
\boxed{
\operatorname{Stretch}_B(x)
=
\operatorname{Max}_{\subseteq}
\{
\Delta_B(x,y):
C(y)=C(x)
\}.
}
\]

The purpose is to destroy false necessities and accidental correlation.

Generation therefore prefers cases that:

\begin{itemize}
\tightlist
\item
  preserve consequence;
\item
  differ strongly in arrangement;
\item
  differ strongly in succession;
\item
  violate the current explanation;
\item
  need not be realistic;
\item
  may be extreme or degenerate.
\end{itemize}

The generation law is:

\[
\boxed{
\textbf{Generate for discrimination, not plausibility.}
}
\]

\begin{center}\rule{0.5\linewidth}{0.5pt}\end{center}

\subsection{15. Squeeze across consequence
classes}\label{squeeze-across-consequence-classes}

The second pressure is:

\[
\boxed{
\text{minimize deformation while consequence changes.}
}
\]

For \(x\in C^b\),

\[
\boxed{
\operatorname{Squeeze}_B(x)
=
\operatorname{Min}_{\subseteq}
\{
\Delta_B(x,y):
C(y)\neq C(x)
\}.
}
\]

More generally, under an arrangement deformation preorder:

\[
\boxed{
\operatorname{Cross}_{C,\rho}(x,b)
=
\operatorname{Front}_{\preceq_\rho^x}
\{
y:
C(y)=b
\land
\operatorname{Adm}_\rho(x,y)
\}.
}
\]

where:

\[
\operatorname{Front}_{\preceq}(S)
=
\{s\in S:\nexists s'\in S,\ s'\prec s\}.
\]

The boundary law is:

\[
\boxed{
\textbf{Between opposite consequences, preserve everything possible
and minimize the change that flips the consequence.}
}
\]

\begin{center}\rule{0.5\linewidth}{0.5pt}\end{center}

\subsection{16. Bidirectional reversal}\label{bidirectional-reversal}

For every positive anchor:

\[
x^+\in C^+,
\]

ask for a minimal reversal into

\[
C^-.
\]

For every negative anchor:

\[
x^-\in C^-,
\]

ask for a minimal reversal into

\[
C^+.
\]

Thus inquiry probes:

\[
\boxed{
C^+\to C^-
}
\]

and

\[
\boxed{
C^-\to C^+.
}
\]

For problems:

\[
P\to\neg P
\]

and:

\[
\neg P\to P.
\]

The boundary is learned from both directions.

\begin{center}\rule{0.5\linewidth}{0.5pt}\end{center}

\section{Part VII --- Factor Discovery}\label{part-vii-factor-discovery}

\subsection{17. Candidate factor}\label{candidate-factor}

Let

\[
h:X\to H
\]

be a candidate representation.

Define its factorization defect:

\[
\boxed{
FD_\Phi(h)
=
\{
(x,y):
h(x)=h(y)
\land
\Phi(x)\neq\Phi(y)
\}.
}
\]

If:

\[
FD_\Phi(h)=\varnothing,
\]

then:

\[
\boxed{
\Phi=d\circ h
}
\]

for some decoder \(d\) on the declared scope.

Representations are ordered by factorization:

\[
\boxed{
h_1\preceq h_2
\iff
\exists m,\ h_1=m\circ h_2.
}
\]

Inquiry seeks a coarsest or nondominated factor that:

\begin{enumerate}
\def\labelenumi{\arabic{enumi}.}
\tightlist
\item
  preserves protected distinctions;
\item
  eliminates unnecessary distinctions;
\item
  remains noncircular.
\end{enumerate}

Abstract interpretation provides the established lattice-theoretic
framework for sound abstraction and approximation, including
abstraction/concretization relations and fixed-point semantics. {[}R4{]}

\begin{center}\rule{0.5\linewidth}{0.5pt}\end{center}

\section{Part VIII --- Mandatory
Falsification}\label{part-viii-mandatory-falsification}

\subsection{18. Necessity}\label{necessity}

Suppose:

\[
C\Rightarrow N.
\]

The corresponding counterexample region is:

\[
\boxed{
C\land\neg N.
}
\]

The claim is attacked by attempting to realize that region.

If realized:

\[
C\not\Rightarrow N.
\]

If independently established empty:

\[
C\land\neg N=\varnothing,
\]

necessity is warranted relative to scope.

Failure to find a counterexample is not itself proof.

\begin{center}\rule{0.5\linewidth}{0.5pt}\end{center}

\subsection{19. Sufficiency}\label{sufficiency}

Suppose:

\[
S\Rightarrow C.
\]

Attack:

\[
\boxed{
S\land\neg C.
}
\]

If realized, sufficiency fails.

If independently established empty, the implication is warranted in
scope.

Thus:

\[
\boxed{
\textbf{every necessity generates its missing counterexample cell;}
}
\]

\[
\boxed{
\textbf{every sufficiency generates its missing counterexample cell.}
}
\]

SAT/SMT backends can execute these obligations directly once the
semantic claims have been reified into formulas. DPLL(T) provides a
modular architecture in which a general Boolean search procedure
cooperates with specialized theory solvers. {[}R5{]}

\begin{center}\rule{0.5\linewidth}{0.5pt}\end{center}

\section{Part IX --- Description Is Not
Control}\label{part-ix-description-is-not-control}

\subsection{20. Descriptive factor}\label{descriptive-factor}

A factor:

\[
\Phi=d\circ h
\]

may perfectly describe consequence without supplying any way to change
it.

Therefore:

\[
\boxed{
\text{descriptive factor}
\neq
\text{control relation}.
}
\]

\begin{center}\rule{0.5\linewidth}{0.5pt}\end{center}

\subsection{21. Control certificate}\label{control-certificate}

Let \(S\subseteq X\) be a source region and \(T\subseteq X\) a target
region.

A succession \(\pi\) is a must-control certificate iff:

\[
\boxed{
S\subseteq\operatorname{MustPre}_\pi(T).
}
\]

For one state:

\[
x\in\operatorname{MustPre}_\pi(T).
\]

When reasoning happens on an abstraction, a separate
concrete-to-abstract control-refinement relation is required. Feedback
refinement relations were developed specifically to guarantee that
symbolic controllers synthesized over abstractions can be concretely
implemented while preserving specifications. {[}R6{]}

\begin{center}\rule{0.5\linewidth}{0.5pt}\end{center}

\section{Part X --- Actualization}\label{part-x-actualization}

\subsection{22. Desired region}\label{desired-region}

Let:

\[
T\subseteq X
\]

be the desired consequence region.

For current arrangement \(x\), define:

\[
\boxed{
\operatorname{Emb}(x,T)
=
\operatorname{Front}_{\preceq_A^x}(T).
}
\]

This seeks target realizations that preserve the present arrangement as
far as possible.

\begin{center}\rule{0.5\linewidth}{0.5pt}\end{center}

\subsection{23. Arrangement obstruction}\label{arrangement-obstruction}

If the target cannot coexist with retained arrangement constraints,
identify a conflict.

Let:

\[
\Gamma_x
\]

be retained constraints and \(I\) the target relation.

A conflict:

\[
K\subseteq\Gamma_x\cup\{I\}
\]

is minimal iff:

\[
K\text{ is inconsistent}
\]

and:

\[
\forall K'\subsetneq K,\quad
K'\text{ is consistent}.
\]

QuickXPlain is an established divide-and-conquer approach for extracting
preferred/minimal explanations or relaxations from overconstrained
systems. {[}R7{]}

\begin{center}\rule{0.5\linewidth}{0.5pt}\end{center}

\subsection{24. Succession obstruction}\label{succession-obstruction}

If compatible target arrangements exist but cannot be reached, define:

\[
\Pi(x,T)
=
\{
\pi:
R_\pi[x]\cap T\neq\varnothing
\}.
\]

If:

\[
\Pi(x,T)=\varnothing,
\]

regress through predecessor conditions.

Any missing predecessor becomes an inquiry obligation.

\begin{center}\rule{0.5\linewidth}{0.5pt}\end{center}

\subsection{25. Prerequisite attack}\label{prerequisite-attack}

Suppose \(N\) is proposed as necessary to reach \(T\).

Define:

\[
Reach_{\neg N}(x)
\]

as the region reachable while keeping \(N\) absent.

Then test:

\[
\boxed{
Reach_{\neg N}(x)\cap T.
}
\]

If nonempty, \(N\) is route-specific.

Only prerequisites that survive alternate-route attack are retained as
genuine requirements.

\begin{center}\rule{0.5\linewidth}{0.5pt}\end{center}

\section{Part XI --- Candidate Claims and
Warrant}\label{part-xi-candidate-claims-and-warrant}

\subsection{26. Claim wrapper}\label{claim-wrapper}

The crucial architecture type is:

\[
\boxed{
\operatorname{Claim}[\tau].
}
\]

A question requesting an object of role \(\tau\) returns a claim of that
role, not an established object.

Formally:

\[
\boxed{
Q_\tau
:
\Sigma
\rightsquigarrow
\operatorname{Claim}[\tau].
}
\]

A claim may contain arbitrary natural-language payload.

Define:

\[
\boxed{
c=
(
id,
role,
args,
payload,
scope,
provenance,
status,
parents,
level
).
}
\]

Where:

\begin{itemize}
\tightlist
\item
  \(id\) uniquely identifies the claim;
\item
  \(role\) is the type requested by the question;
\item
  \(args\) are the bound referents;
\item
  \(payload\) is the answer content;
\item
  \(scope\) records context;
\item
  \(provenance\) records source;
\item
  \(status\) records epistemic standing;
\item
  \(parents\) records dependency;
\item
  \(level\) records formalization depth.
\end{itemize}

\begin{center}\rule{0.5\linewidth}{0.5pt}\end{center}

\subsection{27. Claim status}\label{claim-status}

Recommended statuses:

\[
\boxed{
\{
proposed,
contested,
reified,
supported,
warranted,
refuted,
inactive,
unknown
\}.
}
\]

A claim's question type does not imply truth.

\begin{center}\rule{0.5\linewidth}{0.5pt}\end{center}

\subsection{28. Lazy semantic
formalization}\label{lazy-semantic-formalization}

Semantic formalization proceeds only as far as needed.

\subsubsection{\texorpdfstring{Level \(L_0\): opaque typed
claim}{Level L\_0: opaque typed claim}}\label{level-l_0-opaque-typed-claim}

Stored:

\begin{itemize}
\tightlist
\item
  question role;
\item
  bound references;
\item
  answer payload.
\end{itemize}

No logical parsing required.

\subsubsection{\texorpdfstring{Level \(L_1\): relational
slots}{Level L\_1: relational slots}}\label{level-l_1-relational-slots}

Extract only:

\begin{itemize}
\tightlist
\item
  referenced objects;
\item
  named difference;
\item
  proposed condition;
\item
  transformation;
\item
  scope;
\item
  polarity.
\end{itemize}

\subsubsection{\texorpdfstring{Level \(L_2\): formal
candidate}{Level L\_2: formal candidate}}\label{level-l_2-formal-candidate}

Produce a partial formalization:

\[
Candidate[\tau].
\]

For example:

\begin{itemize}
\tightlist
\item
  predicate;
\item
  relation;
\item
  constraint set;
\item
  transition;
\item
  guard;
\item
  formula;
\item
  path.
\end{itemize}

\subsubsection{\texorpdfstring{Level \(L_3\): warranted
object}{Level L\_3: warranted object}}\label{level-l_3-warranted-object}

A verified or otherwise licensed object enters the hard retained theory.

Thus:

\[
\boxed{
L_0
\to
L_1
\to
L_2
\to
L_3
}
\]

is demand-driven rather than mandatory.

\begin{center}\rule{0.5\linewidth}{0.5pt}\end{center}

\subsection{29. Candidate store}\label{candidate-store}

Let:

\[
\boxed{
\mathcal C_t
}
\]

be the candidate claim graph.

It may contain:

\begin{itemize}
\tightlist
\item
  incompatible claims;
\item
  contradictory claims;
\item
  unverified claims;
\item
  alternate interpretations;
\item
  speculative claims.
\end{itemize}

It is explicitly non-explosive.

From:

\[
P
\]

and:

\[
\neg P
\]

in the candidate store, the system does not infer arbitrary \(Q\).

Instead:

\[
\boxed{
P,\neg P
\to
\operatorname{Conflict}(P,\neg P).
}
\]

Conflict becomes an inquiry obligation.

This architecture is compatible with the spirit of assumption-based
truth maintenance, which was designed to track conclusions under
assumption sets and to work with inconsistent information without
globally retracting everything. {[}R8{]}

\begin{center}\rule{0.5\linewidth}{0.5pt}\end{center}

\section{Part XII --- Warranted Conditional
Theory}\label{part-xii-warranted-conditional-theory}

\subsection{30. Learned lemma}\label{learned-lemma}

A retained learned object is:

\[
\boxed{
\lambda
=
(g,\ell,\sigma,w,\Gamma,\pi)
}
\]

where:

\begin{itemize}
\tightlist
\item
  \(g\) is an activation guard;
\item
  \(\ell\) is the learned relation;
\item
  \(\sigma\) is scope;
\item
  \(w\) is warrant class;
\item
  \(\Gamma\) is dependency set;
\item
  \(\pi\) is proof/certificate/provenance.
\end{itemize}

The permanent store is:

\[
\boxed{
\Theta_t.
}
\]

Learning is append-only:

\[
\boxed{
\Theta_{t+1}
=
\Theta_t\cup\{\lambda\}.
}
\]

This---not the current possibility region---is the primary monotone
ratchet.

\begin{center}\rule{0.5\linewidth}{0.5pt}\end{center}

\subsection{31. Active theory}\label{active-theory}

For current context \(\gamma\):

\[
\boxed{
\Theta_t^\gamma
=
\{
\ell_i:
\gamma\models g_i
\}.
}
\]

The live model region is:

\[
\boxed{
Q_t^\gamma
=
Models(\Theta_t^\gamma).
}
\]

If a guard ceases to hold, its lemma becomes inactive rather than being
erased.

Thus:

\[
\boxed{
\textbf{reopening is guard deactivation, not epistemic amnesia.}
}
\]

\begin{center}\rule{0.5\linewidth}{0.5pt}\end{center}

\subsection{32. Warrant acyclicity}\label{warrant-acyclicity}

Let \(W_t\) be the warrant dependency graph.

For newly accepted warrant \(w_{t+1}\):

\[
\boxed{
Parents(w_{t+1})
\subseteq
V(W_t).
}
\]

Positive warrant may not depend on itself.

Control may recurse.

Warrant may not circle.

\begin{center}\rule{0.5\linewidth}{0.5pt}\end{center}

\section{Part XIII --- Question
Contracts}\label{part-xiii-question-contracts}

\subsection{33. Question contract}\label{question-contract}

Every question family is defined by:

\[
\boxed{
\mathcal Q
=
(
id,
\tau_{in},
\tau_{claim},
Pre,
Render,
Bind,
Reify,
Check,
Update,
Next
).
}
\]

Where:

\begin{itemize}
\tightlist
\item
  \(id\) identifies the question operator;
\item
  \(\tau_{in}\) specifies input referent types;
\item
  \(\tau_{claim}\) specifies the role of the claim produced;
\item
  \(Pre\) determines whether the question is well-typed;
\item
  \(Render\) generates ordinary-language wording;
\item
  \(Bind\) wraps the answer as a typed claim;
\item
  \(Reify\) optionally converts the claim to a formal candidate;
\item
  \(Check\) tests or warrants the candidate;
\item
  \(Update\) changes inquiry state;
\item
  \(Next\) generates possible subsequent obligations.
\end{itemize}

\begin{center}\rule{0.5\linewidth}{0.5pt}\end{center}

\subsection{34. Binding law}\label{binding-law}

Given semantic answer \(a\):

\[
\boxed{
Bind_Q(a)
=
Claim_Q(a).
}
\]

This operation does not assert:

\[
a:\tau.
\]

It asserts only:

\[
\boxed{
\text{“payload }a\text{ was offered as an answer to a question of role }\tau.”}
\]

This operation must be total for ordinary answer payloads.

\begin{center}\rule{0.5\linewidth}{0.5pt}\end{center}

\subsection{35. Reification law}\label{reification-law}

Formalization is partial:

\[
\boxed{
Reify_\tau
:
Claim[\tau]
\rightharpoonup
Candidate[\tau].
}
\]

Failure to reify is itself an inquiry event.

It may generate:

\begin{itemize}
\tightlist
\item
  clarification;
\item
  decomposition;
\item
  re-expression;
\item
  rebind;
\item
  or unknown.
\end{itemize}

\begin{center}\rule{0.5\linewidth}{0.5pt}\end{center}

\subsection{36. Promotion law}\label{promotion-law}

A candidate becomes hard knowledge only through:

\[
\boxed{
Promote(c,\pi)
=
\lambda
}
\]

when:

\[
Check(\pi)=\top.
\]

The central safety rule is:

\[
\boxed{
\textbf{question type may create a claim;
only warrant may create a hard lemma.}
}
\]

\begin{center}\rule{0.5\linewidth}{0.5pt}\end{center}

\section{Part XIV --- Protocol Safety}\label{part-xiv-protocol-safety}

\subsection{37. Protocol-safety
proposition}\label{protocol-safety-proposition}

Assume:

\begin{enumerate}
\def\labelenumi{\arabic{enumi}.}
\tightlist
\item
  every question answer is wrapped as \(\operatorname{Claim}[\tau]\),
  never coerced directly into \(\tau\);
\item
  formal reification is partial and may fail safely;
\item
  candidate-store contradictions do not imply arbitrary claims;
\item
  hard theory updates require independent warrant;
\item
  warrant dependencies are acyclic.
\end{enumerate}

Then arbitrary semantic payloads cannot by themselves make the hard
inquiry state formally unsound.

A bad answer may cause:

\begin{itemize}
\tightlist
\item
  bad candidate claims;
\item
  unproductive branches;
\item
  conflicts;
\item
  extra obligations;
\item
  failed reification;
\item
  or misleading provisional paths.
\end{itemize}

But it cannot lawfully become a hard fact merely by being generated.

This gives the key implementation principle:

\[
\boxed{
\textbf{the question graph must remain lawful even if every answer is wrong.}
}
\]

This does not guarantee useful progress. It guarantees containment of
semantic fallibility.

\begin{center}\rule{0.5\linewidth}{0.5pt}\end{center}

\section{Part XV --- The Adaptive Abstract Transition
System}\label{part-xv-the-adaptive-abstract-transition-system}

\subsection{38. Intermediate structure}\label{intermediate-structure}

The formal architecture maintains:

\[
\boxed{
\mathcal I_t
=
(
B_t,
V_t,
\alpha_t,
\widehat{\mathcal U}_t,
\widehat\Phi_t,
\mathcal C_t,
\Theta_t,
\mathcal O_t,
W_t
).
}
\]

Where:

\begin{itemize}
\tightlist
\item
  \(B_t\): predicate/distinction basis;
\item
  \(V_t\): abstract carrier;
\item
  \(\alpha_t\): abstraction map;
\item
  \(\widehat{\mathcal U}_t\): abstract transitions;
\item
  \(\widehat\Phi_t\): protected abstract consequence;
\item
  \(\mathcal C_t\): candidate claim graph;
\item
  \(\Theta_t\): warranted conditional store;
\item
  \(\mathcal O_t\): unresolved obligations;
\item
  \(W_t\): warrant DAG.
\end{itemize}

\begin{center}\rule{0.5\linewidth}{0.5pt}\end{center}

\subsection{39. Predicate abstraction}\label{predicate-abstraction}

Let:

\[
B_t=\{p_1,\ldots,p_n\}.
\]

Define:

\[
\boxed{
\alpha_t(x)
=
(p_1(x),\ldots,p_n(x)).
}
\]

The realized abstract carrier is:

\[
\boxed{
V_t
=
\operatorname{im}\alpha_t.
}
\]

No nonexistent Boolean combination is silently treated as concrete
reality.

Predicate abstraction and counterexample-guided abstraction refinement
provide established machinery for starting with a coarse abstraction and
adding distinctions when spurious behavior reveals inadequacy. {[}R9{]}

\begin{center}\rule{0.5\linewidth}{0.5pt}\end{center}

\subsection{40. Abstraction/concretization
law}\label{abstractionconcretization-law}

Extend abstraction to sets:

\[
\alpha_t^\sharp(S)
=
\{\alpha_t(x):x\in S\}.
\]

Define concretization:

\[
\gamma_t(A)
=
\{x:\alpha_t(x)\in A\}.
\]

Require:

\[
\boxed{
\alpha_t^\sharp(S)\subseteq A
\iff
S\subseteq\gamma_t(A)
}
\]

when using a Galois-insertion-style exact realized abstraction.

\begin{center}\rule{0.5\linewidth}{0.5pt}\end{center}

\section{Part XVI --- Lawfulness
Conditions}\label{part-xvi-lawfulness-conditions}

\subsection{41. L1 --- Realizability}\label{l1-realizability}

If an abstract state is treated as actual:

\[
v\in V_t
\]

must correspond to at least one concrete realization.

Otherwise its overapproximate status must be explicit.

\begin{center}\rule{0.5\linewidth}{0.5pt}\end{center}

\subsection{42. L2 --- Consequence
soundness}\label{l2-consequence-soundness}

Hard abstract claims must not entail false concrete consequences.

For exact consequence compression:

\[
\boxed{
FD_{\mathcal H}(\alpha_t)=\varnothing.
}
\]

\begin{center}\rule{0.5\linewidth}{0.5pt}\end{center}

\subsection{43. L3 --- Transition
soundness}\label{l3-transition-soundness}

Every concrete transition must be represented abstractly:

\[
\boxed{
xR_uy
\Rightarrow
\alpha_t(x)\widehat R_u\alpha_t(y).
}
\]

\begin{center}\rule{0.5\linewidth}{0.5pt}\end{center}

\subsection{44. L4 --- Control refinement}\label{l4-control-refinement}

Any abstract control action used to guarantee a concrete outcome must
satisfy an implementation/refinement relation sufficient to transport
the guarantee to the concrete system.

\begin{center}\rule{0.5\linewidth}{0.5pt}\end{center}

\subsection{45. L5 --- Conservative
refinement}\label{l5-conservative-refinement}

If \(B_{t+1}\) refines \(B_t\), require a forgetting map:

\[
\boxed{
\rho_{t+1,t}
:
V_{t+1}\to V_t
}
\]

such that:

\[
\boxed{
\alpha_t
=
\rho_{t+1,t}\circ\alpha_{t+1}.
}
\]

Old abstraction is recoverable by forgetting newly introduced
distinctions.

\begin{center}\rule{0.5\linewidth}{0.5pt}\end{center}

\subsection{46. L6 --- Compression
exactness}\label{l6-compression-exactness}

A quotient:

\[
q:V_t\to V'
\]

may be used only if protected consequence factors through it:

\[
\boxed{
\widehat\Phi_t
=
\Phi'\circ q.
}
\]

If control or succession is protected, relevant transition behavior must
factor as well.

\begin{center}\rule{0.5\linewidth}{0.5pt}\end{center}

\subsection{47. L7 --- Hard-lemma
soundness}\label{l7-hard-lemma-soundness}

Every irreversible hard update requires adequate warrant:

\[
\boxed{
Check(\pi)=\top
\Longrightarrow
g\models_\sigma\ell.
}
\]

Proof-producing SMT systems and frameworks such as LFSC illustrate the
established pattern of complex solvers returning certificates that can
be checked by a smaller trusted mechanism. {[}R10{]}

\begin{center}\rule{0.5\linewidth}{0.5pt}\end{center}

\subsection{48. L8 --- Warrant acyclicity}\label{l8-warrant-acyclicity}

The proof graph must remain acyclic.

\begin{center}\rule{0.5\linewidth}{0.5pt}\end{center}

\section{Part XVII --- Established Algorithmic
Spine}\label{part-xvii-established-algorithmic-spine}

\subsection{49. Relational semantics}\label{relational-semantics}

\textbf{Role:} represent succession, composition, predicates, may/must
reachability.

\textbf{Established basis:} relational algebra, predicate transformers,
Kleene algebra with tests.

KAT provides an equational system combining programs/actions with tests.
{[}R1{]}

\begin{center}\rule{0.5\linewidth}{0.5pt}\end{center}

\subsection{50. Abstract interpretation}\label{abstract-interpretation}

\textbf{Role:} lawful abstraction, concretization, consequence-relative
compression, fixed-point approximation.

Cousot and Cousot's abstract interpretation framework explicitly relates
concrete computations to abstract universes and provides the
lattice-theoretic basis for sound approximation. {[}R4{]}

\begin{center}\rule{0.5\linewidth}{0.5pt}\end{center}

\subsection{51. CEGAR}\label{cegar}

\textbf{Role:} split false equivalences when a coarse abstraction
produces spurious behavior.

CEGAR iteratively checks an abstraction and refines it from
counterexamples. {[}R9{]}

RCI correspondence:

\[
\boxed{
\text{protected consequence separator}
\to
\text{new predicate}
\to
\text{refined abstraction}.
}
\]

\begin{center}\rule{0.5\linewidth}{0.5pt}\end{center}

\subsection{\texorpdfstring{52. \(L^*\) and behavioral quotient
learning}{52. L\^{}* and behavioral quotient learning}}\label{l-and-behavioral-quotient-learning}

\textbf{Role:} learn consequence-distinguishable state classes through
queries and counterexamples.

Angluin's \(L^*\) algorithm learns regular languages using membership
and equivalence queries and uses counterexamples to refine the learned
structure. {[}R2{]}

RCI correspondence:

\[
\boxed{
\text{false equivalence}
\to
\text{counterexample}
\to
\text{separator}
\to
\text{new state distinction}.
}
\]

\begin{center}\rule{0.5\linewidth}{0.5pt}\end{center}

\subsection{53. Partition refinement}\label{partition-refinement}

\textbf{Role:} compute coarsest stable partitions when the system is
explicit.

Paige--Tarjan gives efficient algorithms for relational coarsest
partition problems. {[}R3{]}

\begin{center}\rule{0.5\linewidth}{0.5pt}\end{center}

\subsection{54. DPLL(T)}\label{dpllt}

\textbf{Role:} modular backend architecture for satisfiability and
theory-specific obligations.

DPLL(T) provides a general search engine parameterized by theory
solvers, including theory propagation and learned constraints. {[}R5{]}

RCI uses the same architectural principle:

\[
\boxed{
\text{generic inquiry controller}
+
\text{specialized discharge engines}.
}
\]

\begin{center}\rule{0.5\linewidth}{0.5pt}\end{center}

\subsection{55. IC3/PDR}\label{ic3pdr}

\textbf{Role:} property-directed recursion, predecessor obligations,
generalized blocking clauses, inductive accumulation.

Bradley's IC3 incrementally generates inductive clauses relative to
reachability approximations and either strengthens a property to an
invariant or finds a counterexample trace. {[}R11{]}

RCI correspondence:

\[
\boxed{
\text{obligation}
\to
\text{predecessor}
\to
\text{recursive obligation}
}
\]

and:

\[
\boxed{
\text{failure/unreachability}
\to
\text{generalized reusable cut}.
}
\]

\begin{center}\rule{0.5\linewidth}{0.5pt}\end{center}

\subsection{56. Constrained Horn clauses and
Spacer}\label{constrained-horn-clauses-and-spacer}

\textbf{Role:} encode recursively dependent obligations and predecessor
requirements.

Z3's Spacer implements generalized PDR for constrained Horn clauses and
supports infinite-state systems over several SMT theories. {[}R12{]}

An RCI subproblem can therefore compile to a relation such as:

\[
O_j(y)
\leftarrow
\varphi(x,y)
\land
O_i(x).
\]

\begin{center}\rule{0.5\linewidth}{0.5pt}\end{center}

\subsection{57. Craig interpolation}\label{craig-interpolation}

\textbf{Role:} derive a relation over shared vocabulary that summarizes
why two constraint sets cannot coexist.

McMillan demonstrated interpolation as a basis for SAT-based model
checking and abstraction. {[}R13{]}

RCI uses interpolation as one possible implementation of:

\[
\boxed{
\text{conflict}
\to
\text{generalized shared boundary relation}.
}
\]

\begin{center}\rule{0.5\linewidth}{0.5pt}\end{center}

\subsection{58. MUS/MCS and QuickXPlain}\label{musmcs-and-quickxplain}

\textbf{Role:} localize minimal conflict and minimal repair.

QuickXPlain computes preferred explanations and relaxations for
overconstrained systems through a solver-independent divide-and-conquer
method. {[}R7{]}

\begin{center}\rule{0.5\linewidth}{0.5pt}\end{center}

\subsection{59. ATMS}\label{atms}

\textbf{Role:} maintain conclusions under assumption environments,
tolerate incompatible candidate worlds, activate/deactivate
context-dependent claims.

De Kleer's ATMS explicitly represents assumption sets and supports
reasoning with inconsistent information and multiple potential
solutions. {[}R8{]}

\begin{center}\rule{0.5\linewidth}{0.5pt}\end{center}

\subsection{60. Feedback refinement}\label{feedback-refinement}

\textbf{Role:} transport abstract control results to concrete systems.

Feedback refinement relations characterize when abstract symbolic
controllers can be implemented on concrete nondeterministic plants while
preserving specifications. {[}R6{]}

\begin{center}\rule{0.5\linewidth}{0.5pt}\end{center}

\section{Part XVIII --- Mapping RCI to Established
Machinery}\label{part-xviii-mapping-rci-to-established-machinery}

\begin{longtable}[]{@{}
  >{\raggedright\arraybackslash}p{(\columnwidth - 4\tabcolsep) * \real{0.3333}}
  >{\raggedright\arraybackslash}p{(\columnwidth - 4\tabcolsep) * \real{0.3333}}
  >{\raggedright\arraybackslash}p{(\columnwidth - 4\tabcolsep) * \real{0.3333}}@{}}
\toprule\noalign{}
\begin{minipage}[b]{\linewidth}\raggedright
RCI relation
\end{minipage} & \begin{minipage}[b]{\linewidth}\raggedright
Formal object
\end{minipage} & \begin{minipage}[b]{\linewidth}\raggedright
Established machinery
\end{minipage} \\
\midrule\noalign{}
\endhead
\bottomrule\noalign{}
\endlastfoot
succession & relation / program & relational algebra, KAT \\
may predecessor & existential predicate transformer & backward
reachability \\
must predecessor & universal predicate transformer &
weakest-precondition style reasoning \\
consequence quotient & equivalence partition & Myhill--Nerode,
bisimulation, Paige--Tarjan \\
adaptive abstraction & Galois abstraction & abstract interpretation \\
false-equivalence split & new predicate/separator & CEGAR,
interpolation \\
same-class stretching & diverse satisfying assignments & SAT/SMT
enumeration, OMT, MaxSAT \\
cross-class squeezing & minimal correction & MCS, optimization,
QuickXPlain \\
necessity attack & satisfiability of \(C\land\neg N\) & SAT/SMT \\
sufficiency attack & satisfiability of \(N\land\neg C\) & SAT/SMT \\
recursive target descent & Horn clauses & CHC, Spacer \\
generalized failure cut & inductive/blocking lemma & CDCL, PDR \\
conflict explanation & minimal inconsistent subset & MUS, QuickXPlain \\
conditional activation & assumption environment & ATMS \\
control implementation & abstraction refinement relation & feedback
refinement \\
proof checking & certificate & LFSC / proof-producing SMT \\
\end{longtable}

No one existing algorithm implements RCI as a whole.

The novel synthesis lies principally in:

\begin{enumerate}
\def\labelenumi{\arabic{enumi}.}
\tightlist
\item
  making extremal same-class variation and minimal cross-class
  deformation the universal inquiry geometry;
\item
  using plain-language typed question contracts as the interface to
  semantic generators;
\item
  allowing generated answers to remain opaque typed claims rather than
  requiring full semantic parsing;
\item
  treating contradictions created by those claims as recursive inquiry
  obligations;
\item
  unifying all backend returns as guarded reusable relational cuts;
\item
  scheduling questions based on expected consequential structural
  refinement.
\end{enumerate}

\begin{center}\rule{0.5\linewidth}{0.5pt}\end{center}

\section{Part XIX --- Full Inquiry
State}\label{part-xix-full-inquiry-state}

\subsection{61. Runtime state}\label{runtime-state}

The recommended state is:

\[
\boxed{
\Sigma_t
=
(
\mathfrak B,
\gamma_t,
\Phi_t,
B_t,
\alpha_t,
\widehat R_t,
\mathcal C_t,
\Theta_t,
\mathcal O_t,
W_t,
M_t
).
}
\]

Where:

\begin{itemize}
\tightlist
\item
  \(\mathfrak B\): binding;
\item
  \(\gamma_t\): active context;
\item
  \(\Phi_t\): protected consequence profile;
\item
  \(B_t\): active distinction basis;
\item
  \(\alpha_t\): current abstraction;
\item
  \(\widehat R_t\): abstract succession structure;
\item
  \(\mathcal C_t\): candidate claim graph;
\item
  \(\Theta_t\): warranted conditional theory;
\item
  \(\mathcal O_t\): unresolved obligations;
\item
  \(W_t\): warrant DAG;
\item
  \(M_t\): current inquiry mode.
\end{itemize}

\begin{center}\rule{0.5\linewidth}{0.5pt}\end{center}

\section{Part XX --- Obligations}\label{part-xx-obligations}

\subsection{62. Obligation type}\label{obligation-type}

Define:

\[
\boxed{
O
=
(
id,
kind,
carrier,
args,
scope,
priority,
parents
).
}
\]

Canonical kinds include:

\[
\operatorname{Establish}(C),
\]

\[
\operatorname{Eliminate}(C),
\]

\[
\operatorname{Distinguish}(x,y),
\]

\[
\operatorname{TestNec}(N,C),
\]

\[
\operatorname{TestSuff}(N,C),
\]

\[
\operatorname{CrossBoundary}(x,C),
\]

\[
\operatorname{ResolveConflict}(K),
\]

\[
\operatorname{EstablishPrereq}(N),
\]

\[
\operatorname{TestPrereq}(N,T),
\]

\[
\operatorname{LocalizeFailure}(f,s),
\]

\[
\operatorname{Reify}(c),
\]

\[
\operatorname{RefineAbstraction}(x,y),
\]

\[
\operatorname{ValidateWarrant}(c),
\]

\[
\operatorname{Reopen}(g).
\]

\begin{center}\rule{0.5\linewidth}{0.5pt}\end{center}

\section{Part XXI --- Question Modes}\label{part-xxi-question-modes}

\subsection{63. Inquiry modes}\label{inquiry-modes}

Use the following principal modes:

\[
\mathsf V=\text{Variation},
\]

\[
\mathsf B=\text{Boundary},
\]

\[
\mathsf F=\text{Falsification},
\]

\[
\mathsf A=\text{Actualization},
\]

\[
\mathsf W=\text{Warrant},
\]

\[
\mathsf C=\text{Compression}.
\]

A useful regulatory rhythm is:

\[
\boxed{
\mathsf V
\to
\mathsf B
\to
\mathsf F
\to
\mathsf A
\to
\mathsf W
\to
\mathsf C.
}
\]

However, the mature controller is event-driven.

The return determines the next mode.

\begin{center}\rule{0.5\linewidth}{0.5pt}\end{center}

\section{Part XXII --- Main Algorithm}\label{part-xxii-main-algorithm}

\subsection{64. Macro algorithm}\label{macro-algorithm}

\begin{Shaded}
\begin{Highlighting}[]
\NormalTok{RATCHETING CONSEQUENCE INQUIRY}

\NormalTok{INPUT}
\NormalTok{    inquiry binding B}
\NormalTok{    current context γ}
\NormalTok{    protected consequence profile Φ}
\NormalTok{    initial obligation O}
\NormalTok{    semantic generator M}
\NormalTok{    discharge backends E}

\NormalTok{STATE}
\NormalTok{    candidate claim graph C}
\NormalTok{    warranted conditional theory Θ}
\NormalTok{    active predicate basis Bp}
\NormalTok{    abstraction α}
\NormalTok{    abstract transition relation Rhat}
\NormalTok{    obligation set Oset}
\NormalTok{    warrant DAG W}

\NormalTok{LOOP}
\NormalTok{1. SELECT OBLIGATION}
\NormalTok{   Choose an unresolved consequential obligation.}

\NormalTok{2. COMPILE QUESTION}
\NormalTok{   Select a lawful question contract whose preconditions match}
\NormalTok{   the obligation and bound referents.}

\NormalTok{3. RENDER}
\NormalTok{   Render only the locally required plain{-}language question.}

\NormalTok{4. GENERATE}
\NormalTok{   Obtain semantic payload a from the semantic generator.}

\NormalTok{5. BIND}
\NormalTok{   Store:}
\NormalTok{       c := Claim(question\_role, bound\_args, a)}
\NormalTok{   Do not promote c to truth.}

\NormalTok{6. COMPOSE}
\NormalTok{   Relate c to prior claims using the semantics of the question}
\NormalTok{   contracts that produced them.}

\NormalTok{7. DETECT CONSEQUENCES}
\NormalTok{   Determine whether composition creates:}
\NormalTok{       a new distinction,}
\NormalTok{       a counterexample,}
\NormalTok{       a conflict,}
\NormalTok{       a possible equivalence,}
\NormalTok{       a boundary witness,}
\NormalTok{       a prerequisite,}
\NormalTok{       a target,}
\NormalTok{       a reification need,}
\NormalTok{       a warrant need,}
\NormalTok{       or no consequential change.}

\NormalTok{8. IF NEEDED, REIFY}
\NormalTok{   Convert only consequential claims into formal candidate predicates,}
\NormalTok{   relations, constraints, transitions, or guards.}
\NormalTok{   If reification fails, create a reification/decomposition obligation.}

\NormalTok{9. DISCHARGE}
\NormalTok{   Send formalizable obligations to an appropriate backend:}
\NormalTok{       SAT/SMT,}
\NormalTok{       PDR/CHC,}
\NormalTok{       optimization,}
\NormalTok{       conflict extraction,}
\NormalTok{       controller synthesis,}
\NormalTok{       experiment,}
\NormalTok{       observation,}
\NormalTok{       retrieval,}
\NormalTok{       external query,}
\NormalTok{       or another licensed mechanism.}

\NormalTok{10. LOCALIZE}
\NormalTok{    For success:}
\NormalTok{        subtract unnecessary support.}
\NormalTok{    For failure:}
\NormalTok{        isolate arrangement, succession, or interaction residual.}
\NormalTok{    For contradiction:}
\NormalTok{        isolate the smallest scope/guard/meaning distinction that}
\NormalTok{        resolves the conflict.}

\NormalTok{11. ATTACK}
\NormalTok{    Every proposed necessity creates a consequence{-}without{-}it obligation.}
\NormalTok{    Every proposed sufficiency creates a condition{-}without{-}consequence obligation.}
\NormalTok{    Every proposed prerequisite creates a target{-}without{-}prerequisite obligation.}

\NormalTok{12. GENERALIZE}
\NormalTok{    Derive the nondominated frontier of reusable relations warranted}
\NormalTok{    by the return.}

\NormalTok{13. LEARN}
\NormalTok{    Promote only adequately warranted relations into Θ with:}
\NormalTok{        guard,}
\NormalTok{        scope,}
\NormalTok{        provenance,}
\NormalTok{        dependency,}
\NormalTok{        certificate.}

\NormalTok{14. REFINE OR COMPRESS}
\NormalTok{    Split any abstraction that identifies protected{-}consequence{-}distinct}
\NormalTok{    cases.}
\NormalTok{    Merge distinctions no protected consequence or protected control}
\NormalTok{    behavior can inspect.}

\NormalTok{15. REOPEN}
\NormalTok{    If a stored guard no longer holds, reactivate the distinction or}
\NormalTok{    obligation it protected.}

\NormalTok{16. PROGRESS TEST}
\NormalTok{    Retain a semantic step only if it:}
\NormalTok{        adds nonredundant warrant,}
\NormalTok{        eliminates possibility,}
\NormalTok{        splits false equivalence,}
\NormalTok{        justifies equivalence,}
\NormalTok{        localizes a boundary,}
\NormalTok{        falsifies necessity/sufficiency,}
\NormalTok{        resolves a prerequisite,}
\NormalTok{        exposes an obstruction,}
\NormalTok{        or changes the next consequential obligation.}

\NormalTok{17. RECUR}
\NormalTok{    Convert the consequential residual into the next obligation.}

\NormalTok{STOP WHEN}
\NormalTok{    the obligation is warranted satisfied;}
\NormalTok{    the obligation is warranted impossible in the declared binding;}
\NormalTok{    all remaining alternatives are protected{-}consequence{-}equivalent;}
\NormalTok{    or no available question or backend can distinguish the remaining alternatives.}
\end{Highlighting}
\end{Shaded}

\begin{center}\rule{0.5\linewidth}{0.5pt}\end{center}

\section{Part XXIII --- Stretch/Squeeze
Subalgorithm}\label{part-xxiii-stretchsqueeze-subalgorithm}

\subsection{65. Stretch}\label{stretch}

\begin{Shaded}
\begin{Highlighting}[]
\NormalTok{STRETCH(C, explanation E)}

\NormalTok{Generate or retrieve cases satisfying the same consequence C.}

\NormalTok{Prefer cases that:}
\NormalTok{    disagree with E;}
\NormalTok{    differ from previous cases in arrangement;}
\NormalTok{    differ in succession;}
\NormalTok{    remove proposed necessities;}
\NormalTok{    use distinct mechanisms;}
\NormalTok{    are extreme when this improves discrimination.}

\NormalTok{For each surviving common relation N:}
\NormalTok{    create TestNec(N, C).}
\end{Highlighting}
\end{Shaded}

\begin{center}\rule{0.5\linewidth}{0.5pt}\end{center}

\subsection{66. Squeeze}\label{squeeze}

\begin{Shaded}
\begin{Highlighting}[]
\NormalTok{SQUEEZE(x, C)}

\NormalTok{Seek an opposite{-}consequence realization y.}
\NormalTok{Preserve as much of x as possible.}

\NormalTok{Minimize:}
\NormalTok{    arrangement difference,}
\NormalTok{    succession difference,}
\NormalTok{    or a nondominated combination if no common scalar exists.}

\NormalTok{Compare boundary witnesses from both consequence directions.}
\NormalTok{Any residual relation becomes a candidate factor.}
\end{Highlighting}
\end{Shaded}

\begin{center}\rule{0.5\linewidth}{0.5pt}\end{center}

\section{Part XXIV --- Question
Scheduling}\label{part-xxiv-question-scheduling}

\subsection{67. Dependency order}\label{dependency-order}

Define:

\[
q_i\prec q_j
\]

when the answer to \(q_i\) can change:

\begin{itemize}
\tightlist
\item
  the referent of \(q_j\);
\item
  whether \(q_j\) is well-typed;
\item
  what answers to \(q_j\) mean;
\item
  or which outputs are possible.
\end{itemize}

Questions related by \(\prec\) should normally be sequenced.

Independent questions may be batched.

Thus:

\[
\boxed{
\textbf{batch antichains; sequence dependencies.}
}
\]

\begin{center}\rule{0.5\linewidth}{0.5pt}\end{center}

\subsection{68. Order robustness}\label{order-robustness}

For questions expected to be independent, compare:

\[
U_{q_2}(U_{q_1}(\Sigma))
\]

with:

\[
U_{q_1}(U_{q_2}(\Sigma)).
\]

If protected consequences differ, inquire whether:

\begin{enumerate}
\def\labelenumi{\arabic{enumi}.}
\tightlist
\item
  the underlying relations genuinely interact; or
\item
  the language model exhibits irrelevant order sensitivity.
\end{enumerate}

Question-order permutation is therefore itself a diagnostic.

\begin{center}\rule{0.5\linewidth}{0.5pt}\end{center}

\subsection{69. Question value}\label{question-value}

Without a justified scalar objective, retain nondominated questions.

Possible dimensions:

\begin{itemize}
\tightlist
\item
  cost;
\item
  expected discrimination;
\item
  expected generalizability;
\item
  risk;
\item
  destructiveness;
\item
  reversibility;
\item
  warrant quality.
\end{itemize}

With a valid measure \(\nu\):

\[
q^\star
\in
\arg\min_q
\left[
Cost(q)
+
\lambda
\max_r
\nu
(
Q_t\cap[\![Gen(q,r)]\!]
)
\right].
\]

This is optional policy, not foundational structure.

\begin{center}\rule{0.5\linewidth}{0.5pt}\end{center}

\section{Part XXV --- Full Typed Question
Library}\label{part-xxv-full-typed-question-library}

The following questions are intentionally domain-agnostic and
goal-agnostic.

The surface wording may vary. Their formal roles may not.

\begin{center}\rule{0.5\linewidth}{0.5pt}\end{center}

\subsection{70. Present-obligation
questions}\label{present-obligation-questions}

\textbf{Input:} inquiry state\\
\textbf{Output:} \texttt{Claim{[}ObligationCharacterization{]}}

\begin{enumerate}
\def\labelenumi{\arabic{enumi}.}
\tightlist
\item
  What, exactly, is unresolved here?
\item
  Is the unresolved thing something to establish, eliminate,
  distinguish, compare, explain, actualize, repair, test, or determine
  possible?
\item
  What kind of thing could constitute an answer to this?
\item
  Which parts of the present situation are already fixed?
\item
  Which parts remain open?
\item
  Which possible answers would lead to different consequences?
\item
  Which possible answers would lead to different next questions?
\item
  What would count as an independently constrained return?
\end{enumerate}

\begin{center}\rule{0.5\linewidth}{0.5pt}\end{center}

\subsection{71. Binding questions}\label{binding-questions}

\textbf{Output roles:} \texttt{Claim{[}Carrier{]}},
\texttt{Claim{[}Relation{]}}, \texttt{Claim{[}Operation{]}},
\texttt{Claim{[}ExpressibilityGap{]}}

\begin{enumerate}
\def\labelenumi{\arabic{enumi}.}
\tightlist
\item
  What kinds of things can this inquiry currently refer to?
\item
  What type of thing is the present object?
\item
  What distinctions can currently be expressed about it?
\item
  What relations can currently be expressed?
\item
  What changes to it are admissible to consider?
\item
  What transformations are admissible to consider?
\item
  Which consequences are currently protected?
\item
  What available mechanisms could independently constrain an answer?
\item
  Am I assuming a distance that has not been supplied?
\item
  Am I assuming probability that has not been supplied?
\item
  Am I assuming causality that has not been supplied?
\item
  Am I assuming temporal order that has not been supplied?
\item
  Am I assuming determinism that has not been supplied?
\item
  Is the distinction I need expressible in the present representation?
\item
  If not, what new relation or representation would make it expressible?
\item
  Can the old representation be recovered merely by forgetting that new
  distinction?
\end{enumerate}

\begin{center}\rule{0.5\linewidth}{0.5pt}\end{center}

\subsection{72. Consequence questions}\label{consequence-questions}

\textbf{Output:} \texttt{Claim{[}ProtectedConsequence{]}}

\begin{enumerate}
\def\labelenumi{\arabic{enumi}.}
\tightlist
\item
  With respect to what consequence are these possibilities being
  compared?
\item
  What would count as one side of that consequence distinction?
\item
  What would count as its opposite?
\item
  If something is a problem, what would it mean merely for that problem
  not to be present?
\item
  Would every non-problem state count as acceptable?
\item
  If not, what additional consequence separates acceptable from
  unacceptable non-problem states?
\item
  If no goal is given, what distinction still changes something
  protected?
\item
  If only relative improvement is known, what makes one possibility
  preferable to another?
\item
  What would make this consequence trivially true?
\item
  What would make its opposite trivially true?
\end{enumerate}

\begin{center}\rule{0.5\linewidth}{0.5pt}\end{center}

\subsection{73. Distinction questions}\label{distinction-questions}

\textbf{Output:} \texttt{Claim{[}Separator{]}}

\begin{enumerate}
\def\labelenumi{\arabic{enumi}.}
\tightlist
\item
  What are the two things being distinguished?
\item
  Under what one relation are they comparable?
\item
  What does that relation return differently?
\item
  Would they still differ if that relation were removed from
  consideration?
\item
  Does their bare nonidentity matter?
\item
  What would have to remain fixed for this to remain the same
  distinction?
\end{enumerate}

\begin{center}\rule{0.5\linewidth}{0.5pt}\end{center}

\subsection{74. Equivalence questions}\label{equivalence-questions}

\textbf{Output:} \texttt{Claim{[}Separator{]}} or
\texttt{Claim{[}Equivalence{]}}

\begin{enumerate}
\def\labelenumi{\arabic{enumi}.}
\tightlist
\item
  Can any protected continuation distinguish these two?
\item
  If so, which one?
\item
  What consequence differs?
\item
  If none separates them, what reason remains to distinguish them
  actively?
\item
  Can they be treated as one state for this inquiry?
\item
  What would make their suppressed difference consequential again?
\item
  Has a newly admitted consequence split a previously folded class?
\item
  Are these truly equivalent, or have we merely failed to distinguish
  them so far?
\end{enumerate}

\begin{center}\rule{0.5\linewidth}{0.5pt}\end{center}

\subsection{75. Second-distinction / tetrad
questions}\label{second-distinction-tetrad-questions}

\textbf{Output:} \texttt{Claim{[}Predicate{]}},
\texttt{Claim{[}CellWitness{]}}

\begin{enumerate}
\def\labelenumi{\arabic{enumi}.}
\tightlist
\item
  What second distinction could explain why these sides have different
  consequences?
\item
  Can that distinction be present while the consequence is present?
\item
  Can it be absent while the consequence is present?
\item
  Can it be present while the consequence is absent?
\item
  Can both it and the consequence be absent?
\item
  Which combinations are expressible?
\item
  Which are constructible?
\item
  Which are compatible?
\item
  Which are realizable?
\item
  Which are reachable?
\item
  Which are actual?
\item
  If a cell is missing, is it impossible or merely unobserved?
\item
  Which missing cell would refute necessity?
\item
  Which missing cell would refute sufficiency?
\end{enumerate}

\begin{center}\rule{0.5\linewidth}{0.5pt}\end{center}

\subsection{76. Arrangement questions}\label{arrangement-questions}

\textbf{Output:} \texttt{Claim{[}ArrangementDifference{]}}

\begin{enumerate}
\def\labelenumi{\arabic{enumi}.}
\tightlist
\item
  What differs about what is present?
\item
  What differs about what is absent?
\item
  What differs about how the parts are related?
\item
  What changes if the arrangement changes while the path is held fixed?
\item
  Which arrangement change is smaller relative to the reference state?
\item
  Can this deformation be weakened while preserving the same
  consequence?
\end{enumerate}

\begin{center}\rule{0.5\linewidth}{0.5pt}\end{center}

\subsection{77. Succession questions}\label{succession-questions}

\textbf{Output:} \texttt{Claim{[}SuccessionDifference{]}}

\begin{enumerate}
\def\labelenumi{\arabic{enumi}.}
\tightlist
\item
  What differs only about how one state becomes another?
\item
  What differs in order?
\item
  What differs in dependency?
\item
  What differs in path?
\item
  What changes if the path changes while the relevant arrangement is
  held fixed?
\item
  Can this succession be weakened while preserving the result?
\item
  If neither arrangement nor succession alone explains the difference,
  does their interaction do so?
\end{enumerate}

\begin{center}\rule{0.5\linewidth}{0.5pt}\end{center}

\subsection{78. Stretch-positive
questions}\label{stretch-positive-questions}

\textbf{Output:} \texttt{Claim{[}SameClassWitness{]}}

\begin{enumerate}
\def\labelenumi{\arabic{enumi}.}
\tightlist
\item
  What could make this consequence unmistakably true while sharing as
  little as possible with the cases already considered?
\item
  What radically different arrangement preserves the same consequence?
\item
  What radically different succession preserves it?
\item
  What extreme or implausible case would preserve the consequence while
  violating as much of the current explanation as possible?
\item
  Can the consequence occur without this proposed feature?
\item
  Can it occur through a completely different mechanism?
\item
  What else can vary while it remains true?
\item
  Can the same-consequence cases differ along one more live distinction?
\item
  What still survives after maximally different positive cases are
  compared?
\item
  Has that survivor actually been shown necessary?
\end{enumerate}

\begin{center}\rule{0.5\linewidth}{0.5pt}\end{center}

\subsection{79. Stretch-negative
questions}\label{stretch-negative-questions}

\textbf{Output:} \texttt{Claim{[}SameClassWitness{]}}

\begin{enumerate}
\def\labelenumi{\arabic{enumi}.}
\tightlist
\item
  What could make the opposite consequence unmistakably true while
  sharing as little as possible with existing negative cases?
\item
  What radically different arrangement gives the same negative
  consequence?
\item
  What radically different path gives the same negative consequence?
\item
  What shared structure is present only because the negative cases have
  not varied enough?
\item
  What proposed explanation of failure can be violated while preserving
  failure?
\end{enumerate}

\begin{center}\rule{0.5\linewidth}{0.5pt}\end{center}

\subsection{80. Squeeze-positive-to-negative
questions}\label{squeeze-positive-to-negative-questions}

\textbf{Output:} \texttt{Claim{[}BoundaryWitness{]}}

\begin{enumerate}
\def\labelenumi{\arabic{enumi}.}
\tightlist
\item
  Starting from this positive case, what is the least arrangement change
  that could make the consequence negative?
\item
  Can more of the original arrangement be preserved while still flipping
  it?
\item
  What is the least succession change that could flip it?
\item
  Can that succession be weakened further?
\end{enumerate}

\begin{center}\rule{0.5\linewidth}{0.5pt}\end{center}

\subsection{81. Squeeze-negative-to-positive
questions}\label{squeeze-negative-to-positive-questions}

\textbf{Output:} \texttt{Claim{[}BoundaryWitness{]}}

\begin{enumerate}
\def\labelenumi{\arabic{enumi}.}
\tightlist
\item
  Starting from this negative case, what is the least arrangement change
  that could make the consequence positive?
\item
  Can more of the original arrangement be preserved?
\item
  What is the least succession change that could make the consequence
  positive?
\item
  Can that succession be weakened further?
\item
  What relation is common to minimal crossings from both directions?
\item
  If the directions differ, what explains the asymmetry?
\end{enumerate}

\begin{center}\rule{0.5\linewidth}{0.5pt}\end{center}

\subsection{82. Extreme reversal
questions}\label{extreme-reversal-questions}

\textbf{Output:} \texttt{Claim{[}Breaker{]}},
\texttt{Claim{[}HiddenGuard{]}}

\begin{enumerate}
\def\labelenumi{\arabic{enumi}.}
\tightlist
\item
  If this makes the consequence overwhelmingly obvious, what would
  nevertheless have to be true for it to fail?
\item
  What is the least additional difference required to make that reversal
  coherent?
\item
  Can that reversing condition be weakened?
\item
  If this makes the opposite overwhelmingly obvious, what would
  nevertheless make the original consequence hold?
\item
  What least additional difference makes that reversal coherent?
\item
  Does this reveal an unstated guard?
\item
  Does this reveal that the consequence itself was defined too coarsely?
\end{enumerate}

\begin{center}\rule{0.5\linewidth}{0.5pt}\end{center}

\subsection{83. Problem/non-problem
questions}\label{problemnon-problem-questions}

\textbf{Output:} \texttt{Claim{[}ProblemBoundary{]}}

\begin{enumerate}
\def\labelenumi{\arabic{enumi}.}
\tightlist
\item
  If the problem is present, what is the least change that would make it
  absent?
\item
  If the problem were absent, what is the least change that would make
  it present?
\item
  What can be removed while the problem remains?
\item
  What must be introduced before the problem appears?
\item
  What relation appears in both problem-removal and
  problem-reconstruction?
\item
  If the two directions differ, what explains the difference?
\item
  Once the problem is absent, what other protected consequence could
  still distinguish this from an acceptable resolution?
\end{enumerate}

\begin{center}\rule{0.5\linewidth}{0.5pt}\end{center}

\subsection{84. Factor questions}\label{factor-questions}

\textbf{Output:} \texttt{Claim{[}Factor{]}}

\begin{enumerate}
\def\labelenumi{\arabic{enumi}.}
\tightlist
\item
  What smaller relation appears to determine the protected consequence?
\item
  Can two cases be identical under that relation while differing in
  consequence?
\item
  If so, what missing distinction separates them?
\item
  Can that distinction be added without changing anything else?
\item
  Can the factor be made coarser while preserving consequence?
\item
  What represented distinction can be forgotten?
\item
  Is the proposed factor merely a renaming of the consequence?
\item
  Can the factor be expressed without invoking the consequence itself?
\item
  Do several nondominated factors remain?
\item
  What question would separate them?
\end{enumerate}

\begin{center}\rule{0.5\linewidth}{0.5pt}\end{center}

\subsection{85. Necessity questions}\label{necessity-questions}

\textbf{Output:} \texttt{Claim{[}NecessityCounterexample{]}} or
\texttt{Claim{[}NoCounterexampleCertificate{]}}

\begin{enumerate}
\def\labelenumi{\arabic{enumi}.}
\tightlist
\item
  Can the consequence still hold if this is false?
\item
  Can this be removed from a positive case while preserving everything
  else possible?
\item
  If removing it breaks the consequence, can something else compensate?
\item
  Can the remaining structure be rearranged so the consequence returns
  without restoring it?
\item
  If that works, what role did the original condition actually play?
\item
  If every attempt fails, what common reason causes those failures?
\item
  Is that common reason a deeper necessity?
\item
  Can that deeper necessity itself be removed?
\item
  Was the search exhaustive?
\item
  What independent result would establish that the
  consequence-without-this region is empty?
\end{enumerate}

\begin{center}\rule{0.5\linewidth}{0.5pt}\end{center}

\subsection{86. Sufficiency questions}\label{sufficiency-questions}

\textbf{Output:} \texttt{Claim{[}SufficiencyCounterexample{]}} or
\texttt{Claim{[}NoCounterexampleCertificate{]}}

\begin{enumerate}
\def\labelenumi{\arabic{enumi}.}
\tightlist
\item
  Can this be true while the consequence is false?
\item
  Can this condition be made strongly present while the consequence
  still fails?
\item
  What additional condition had to change?
\item
  Does that reveal an unstated guard?
\item
  Can the consequence remain without that guard?
\item
  Has the counterexample region been established empty, or merely not
  encountered?
\item
  What independent return would warrant sufficiency?
\end{enumerate}

\begin{center}\rule{0.5\linewidth}{0.5pt}\end{center}

\subsection{87. Abstraction questions}\label{abstraction-questions}

\textbf{Output:} \texttt{Claim{[}Separator{]}},
\texttt{Claim{[}Equivalence{]}}, \texttt{Claim{[}AbstractionDefect{]}}

\begin{enumerate}
\def\labelenumi{\arabic{enumi}.}
\tightlist
\item
  Does every abstract state treated as realized correspond to a concrete
  realization?
\item
  If not, is its overapproximate status explicit?
\item
  Can two concrete realizations map to the same abstract state while
  differing in protected consequence?
\item
  What predicate would separate them?
\item
  Does every concrete transition have an abstract representative?
\item
  Can the old abstraction be recovered by forgetting the new
  distinction?
\item
  Can old warranted statements be transported into the refined
  representation?
\item
  Can this abstraction be compressed further without losing protected
  consequence?
\item
  Can it be compressed without losing protected succession or control
  behavior?
\item
  What first distinction would be lost by one more compression step?
\end{enumerate}

\begin{center}\rule{0.5\linewidth}{0.5pt}\end{center}

\subsection{88. Description/control
questions}\label{descriptioncontrol-questions}

\textbf{Output:} \texttt{Claim{[}ControlRelation{]}}

\begin{enumerate}
\def\labelenumi{\arabic{enumi}.}
\tightlist
\item
  Does this relation merely predict the consequence, or can it be
  established?
\item
  What admissible transformation would establish it?
\item
  Is that transformation available here?
\item
  Can it be implemented from every concrete state represented by this
  abstract state?
\item
  Could any allowed successor still leave the target region?
\item
  Does the proposed route merely permit success, or guarantee it?
\item
  What connects the abstract action to a concrete action?
\item
  Does every concrete successor remain represented abstractly?
\item
  What additional distinction is needed if not?
\item
  What would constitute a genuine control certificate?
\end{enumerate}

\begin{center}\rule{0.5\linewidth}{0.5pt}\end{center}

\subsection{89. Actualization questions}\label{actualization-questions}

\textbf{Output:} \texttt{Claim{[}Target{]}},
\texttt{Claim{[}Conflict{]}}, \texttt{Claim{[}Path{]}}

\begin{enumerate}
\def\labelenumi{\arabic{enumi}.}
\tightlist
\item
  What consequence region counts as the desired result for this
  obligation?
\item
  What target realization preserves the most current structure?
\item
  Can the target coexist with the present arrangement?
\item
  If not, what smallest set of retained conditions makes coexistence
  impossible?
\item
  What can be changed to dissolve that conflict?
\item
  Are several repairs incomparable?
\item
  What distinguishes which repair should be preferred?
\item
  If coexistence is possible, can the target be reached?
\item
  What succession could reach it?
\item
  Does that succession permit or guarantee the target?
\item
  Is there a nondominated alternative?
\end{enumerate}

\begin{center}\rule{0.5\linewidth}{0.5pt}\end{center}

\subsection{90. Prerequisite questions}\label{prerequisite-questions}

\textbf{Output:} \texttt{Claim{[}Prerequisite{]}}

\begin{enumerate}
\def\labelenumi{\arabic{enumi}.}
\tightlist
\item
  What must already be true before this transformation can establish the
  target?
\item
  Which part of that is already true?
\item
  Which part is missing?
\item
  What would have to be true for the missing part itself to hold?
\item
  Is that the first unresolved subproblem?
\item
  Can the target be reached without this proposed prerequisite?
\item
  Can another succession bypass it?
\item
  Is it genuinely necessary or merely route-specific?
\item
  If every route without it fails, what common relation explains the
  failure?
\item
  Can that common relation itself be falsified as necessary?
\end{enumerate}

\begin{center}\rule{0.5\linewidth}{0.5pt}\end{center}

\subsection{91. Discharge-selection
questions}\label{discharge-selection-questions}

\textbf{Output:} \texttt{Claim{[}DischargeMethod{]}}

\begin{enumerate}
\def\labelenumi{\arabic{enumi}.}
\tightlist
\item
  What exact proposition or relation is being tested?
\item
  What kind of return would settle it?
\item
  Would a witness settle it?
\item
  Would an impossibility proof settle it?
\item
  Would a predecessor settle it?
\item
  Would a minimal conflict settle it?
\item
  Would observation settle it?
\item
  Would execution settle it?
\item
  Would retrieval settle it?
\item
  Which available mechanism can produce the needed return with adequate
  warrant?
\item
  What exact query must be sent?
\item
  What raw return did it provide?
\item
  What interpretation does the return license?
\item
  What interpretation does it not license?
\end{enumerate}

\begin{center}\rule{0.5\linewidth}{0.5pt}\end{center}

\subsection{92. Warrant questions}\label{warrant-questions}

\textbf{Output:} \texttt{Claim{[}Warrant{]}}

\begin{enumerate}
\def\labelenumi{\arabic{enumi}.}
\tightlist
\item
  What independently checks this return?
\item
  What exactly would a successful check establish?
\item
  Under what assumptions?
\item
  Over what scope?
\item
  Is the return a proof, observation, execution, retrieval, simulation,
  hypothesis, or another warrant class?
\item
  Does it support a hard exclusion or only provisional preference?
\item
  What stronger statement would exceed the warrant?
\item
  Does the warrant depend on its own conclusion?
\item
  What independent antecedent would remove that circularity?
\end{enumerate}

\begin{center}\rule{0.5\linewidth}{0.5pt}\end{center}

\subsection{93. Success-localization
questions}\label{success-localization-questions}

\textbf{Output:} \texttt{Claim{[}MinimalSupport{]}}

\begin{enumerate}
\def\labelenumi{\arabic{enumi}.}
\tightlist
\item
  What was present in the successful case?
\item
  What can be removed while success remains?
\item
  What can be changed?
\item
  What can be replaced?
\item
  What can be reordered?
\item
  Which survivor has not yet been adequately challenged?
\item
  Can success be reconstructed without it?
\item
  What is the smallest support still sufficient here?
\item
  Is that support sufficient only here or across scope?
\item
  What counterexample would distinguish those possibilities?
\end{enumerate}

\begin{center}\rule{0.5\linewidth}{0.5pt}\end{center}

\subsection{94. Failure-localization
questions}\label{failure-localization-questions}

\textbf{Output:} \texttt{Claim{[}FailureResidual{]}}

\begin{enumerate}
\def\labelenumi{\arabic{enumi}.}
\tightlist
\item
  What is the nearest relevant successful case?
\item
  What differs in arrangement?
\item
  What differs in succession?
\item
  If the arrangement difference is removed, does failure remain?
\item
  If the succession difference is removed, does failure remain?
\item
  If neither alone suffices, does their interaction?
\item
  What part of the failure difference can be removed while failure
  remains?
\item
  Can the residual be reduced further?
\item
  What larger class of failures follows from that residual?
\item
  Under what condition does the generalized failure relation stop
  applying?
\end{enumerate}

\begin{center}\rule{0.5\linewidth}{0.5pt}\end{center}

\subsection{95. Contradiction-repair
questions}\label{contradiction-repair-questions}

\textbf{Output:} \texttt{Claim{[}ConflictResolution{]}}

\begin{enumerate}
\def\labelenumi{\arabic{enumi}.}
\tightlist
\item
  Can these two commitments both hold here?
\item
  If not, which one fails?
\item
  Is one claim too broad?
\item
  Do they apply under different scopes?
\item
  Did a term change meaning?
\item
  Is there a hidden guard separating them?
\item
  Is the supposed counterexample invalid?
\item
  What is the smallest distinction that makes the two claims compatible?
\item
  Does that distinction actually apply here?
\item
  Can the conflict still occur without that distinction?
\item
  Which earlier commitment must be restricted, refuted, or reopened?
\end{enumerate}

\begin{center}\rule{0.5\linewidth}{0.5pt}\end{center}

\subsection{96. Generalization
questions}\label{generalization-questions}

\textbf{Output:} \texttt{Claim{[}GeneralizedRelation{]}}

\begin{enumerate}
\def\labelenumi{\arabic{enumi}.}
\tightlist
\item
  What broader relation is supported by the same reason as this return?
\item
  Can its conditions be weakened?
\item
  Can its conclusion be strengthened?
\item
  Are several incomparable generalizations supported?
\item
  Which are dominated?
\item
  What counterexample would falsify this generalization?
\item
  Does it survive that attack?
\item
  What is the broadest reusable relation that remains warranted?
\end{enumerate}

\begin{center}\rule{0.5\linewidth}{0.5pt}\end{center}

\subsection{97. Conditional-learning
questions}\label{conditional-learning-questions}

\textbf{Output:} \texttt{Claim{[}Guard{]}}, \texttt{Claim{[}Scope{]}}

\begin{enumerate}
\def\labelenumi{\arabic{enumi}.}
\tightlist
\item
  Under what condition is this learned relation active?
\item
  Over what scope?
\item
  Which earlier results does it depend on?
\item
  What warrant supports it?
\item
  What does it let us exclude?
\item
  What does it let us merge?
\item
  What does it force us to split?
\item
  Is it already implied by the retained theory?
\item
  If so, does storing it add anything?
\item
  What condition would make it stop applying?
\item
  Would that make it false or merely inactive?
\end{enumerate}

\begin{center}\rule{0.5\linewidth}{0.5pt}\end{center}

\subsection{98. Reopening questions}\label{reopening-questions}

\textbf{Output:} \texttt{Claim{[}ReopenCondition{]}}

\begin{enumerate}
\def\labelenumi{\arabic{enumi}.}
\tightlist
\item
  Which folded distinction was suppressed under a condition that no
  longer holds?
\item
  What changed?
\item
  Which stored relations depended on the old condition?
\item
  Which become inactive?
\item
  Which remain warranted?
\item
  What alternatives must now be represented separately?
\item
  Can the old relation remain stored under its original guard?
\end{enumerate}

\begin{center}\rule{0.5\linewidth}{0.5pt}\end{center}

\subsection{99. Compression questions}\label{compression-questions}

\textbf{Output:} \texttt{Claim{[}Equivalence{]}}

\begin{enumerate}
\def\labelenumi{\arabic{enumi}.}
\tightlist
\item
  Which active distinctions can no protected consequence inspect?
\item
  Can they be merged without changing protected consequence?
\item
  Can they be merged without changing protected succession behavior?
\item
  Can they be merged without invalidating control guarantees?
\item
  What is the coarsest representation that still answers all protected
  questions correctly?
\item
  What further merge would first create an error?
\item
  What new separator requires a split?
\item
  Can that separator be added conservatively?
\end{enumerate}

\begin{center}\rule{0.5\linewidth}{0.5pt}\end{center}

\subsection{100. Intermediate-lawfulness
questions}\label{intermediate-lawfulness-questions}

\textbf{Output:} \texttt{Claim{[}ArchitectureDefect{]}}

\begin{enumerate}
\def\labelenumi{\arabic{enumi}.}
\tightlist
\item
  Does every abstract state treated as real correspond to a concrete
  state?
\item
  Can any hard abstract consequence be false concretely?
\item
  Can any concrete transition occur without abstract representation?
\item
  Can an abstract control action fail to be concretely implementable?
\item
  Can refinement recover the old abstraction by forgetting?
\item
  Does compression preserve all protected consequences?
\item
  Does it preserve protected transition behavior?
\item
  Does every hard lemma have adequate warrant?
\item
  Does any warrant depend positively on itself?
\item
  What smallest architecture change would restore the violated law?
\end{enumerate}

\begin{center}\rule{0.5\linewidth}{0.5pt}\end{center}

\subsection{101. Question-selection
questions}\label{question-selection-questions}

\textbf{Output:} \texttt{Claim{[}NextQuestionPreference{]}}

\begin{enumerate}
\def\labelenumi{\arabic{enumi}.}
\tightlist
\item
  Which unresolved distinction could still change a protected
  consequence?
\item
  Which could change the next obligation?
\item
  Which available question separates the largest consequentially
  different region?
\item
  Which is likely to yield the broadest reusable relation?
\item
  What would each possible answer let us eliminate, split, merge,
  establish, or reopen?
\item
  Is another question cheaper, safer, more reversible, or better
  warranted without producing a weaker result?
\item
  If no question dominates, which remain nondominated?
\item
  Do all answers to this proposed question lead to the same
  consequential state?
\item
  If so, why ask it?
\end{enumerate}

\begin{center}\rule{0.5\linewidth}{0.5pt}\end{center}

\subsection{102. Progress questions}\label{progress-questions}

\textbf{Output:} \texttt{Claim{[}ProgressStatus{]}}

\begin{enumerate}
\def\labelenumi{\arabic{enumi}.}
\tightlist
\item
  Did this establish a new nonredundant warranted relation?
\item
  Did it remove a live possibility?
\item
  Did it falsify a necessity?
\item
  Did it falsify a sufficiency?
\item
  Did it expose a separator?
\item
  Did it justify an equivalence?
\item
  Did it localize a boundary?
\item
  Did it discharge a prerequisite?
\item
  Did it identify an obstruction?
\item
  Did it establish a reopening condition?
\item
  Did it change any downstream obligation?
\item
  If none occurred, what semantic progress was made?
\item
  If none was made, why should this step enter retained history?
\end{enumerate}

\begin{center}\rule{0.5\linewidth}{0.5pt}\end{center}

\subsection{103. Residual questions}\label{residual-questions}

\textbf{Output:} \texttt{Claim{[}NextObligation{]}}

\begin{enumerate}
\def\labelenumi{\arabic{enumi}.}
\tightlist
\item
  What consequential difference remains unresolved?
\item
  Is it a counterexample?
\item
  Is it an unexplained boundary?
\item
  Is it an arrangement conflict?
\item
  Is it a succession obstruction?
\item
  Is it an unverified prerequisite?
\item
  Is it a factorization defect?
\item
  Is it an expressibility failure?
\item
  Is it a warrant gap?
\item
  What consequence distinction does that residual create?
\item
  What are its two sides?
\item
  Under what relation are they comparable?
\item
  What second distinction could explain them?
\item
  What question now separates them?
\end{enumerate}

\begin{center}\rule{0.5\linewidth}{0.5pt}\end{center}

\subsection{104. Stopping questions}\label{stopping-questions}

\textbf{Output:} \texttt{Claim{[}StopStatus{]}}

\begin{enumerate}
\def\labelenumi{\arabic{enumi}.}
\tightlist
\item
  Is the present obligation warranted satisfied?
\item
  Is it warranted impossible within the declared binding?
\item
  Do all remaining possibilities have the same protected consequences?
\item
  Can any available continuation distinguish them?
\item
  Can any available question change the next warranted action or
  conclusion?
\item
  If not, what consequential reason remains to continue?
\item
  If a distinction remains but cannot be represented, asked, observed,
  or discharged, is the proper result unknown under the present binding?
\end{enumerate}

\begin{center}\rule{0.5\linewidth}{0.5pt}\end{center}

\section{Part XXVI --- Minimal Question
Kernel}\label{part-xxvi-minimal-question-kernel}

\subsection{105. Four primary question
families}\label{four-primary-question-families}

Most of the library can be regenerated from four primary pressures.

\subsubsection{Variation}\label{variation}

\[
\boxed{
\text{What can change while this consequence remains the same?}
}
\]

\subsubsection{Boundary}\label{boundary}

\[
\boxed{
\text{What is the least change that makes the consequence different?}
}
\]

\subsubsection{Falsification}\label{falsification}

\[
\boxed{
\text{Can the consequence still occur without this?}
}
\]

and:

\[
\boxed{
\text{Can this still occur without the consequence?}
}
\]

\subsubsection{Actualization}\label{actualization}

\[
\boxed{
\text{What would have to be true for this relation to hold here?}
}
\]

followed by:

\[
\boxed{
\text{Can it hold without that supposed requirement?}
}
\]

Two additional questions complete the epistemic ratchet:

\[
\boxed{
\text{What does this return warrant beyond this single case?}
}
\]

and:

\[
\boxed{
\text{What distinction can now be forgotten, and what would make it matter again?}
}
\]

\begin{center}\rule{0.5\linewidth}{0.5pt}\end{center}

\section{Part XXVII --- Question-Only Executable
Form}\label{part-xxvii-question-only-executable-form}

The entire architecture may be driven by the following recursive
sequence:

What is unresolved?

What consequence makes the live alternatives different?

What would make that consequence unmistakably true?

What would make its opposite unmistakably true?

What can be radically different while the consequence stays the same?

What can differ as little as possible while making the consequence flip?

What relation survives those comparisons?

Can the consequence still hold without that relation?

Can that relation still hold while the consequence fails?

What smaller relation survives both attacks?

Can any protected consequence distinguish cases that relation
identifies?

If so, what missing distinction separates them?

If not, what distinction can now be ignored?

Does the surviving relation merely describe the consequence, or can it
be established?

What would have to be true for it to hold here?

Can it coexist with what is already here?

If not, what smallest conflict prevents it?

What smallest change dissolves that conflict?

If it can coexist, can it be reached?

If not, what would have to be true beforehand?

Can it be reached without that proposed prerequisite?

If not, what would establish the prerequisite?

What independently constrained mechanism can answer the remaining
question?

What did that mechanism actually return?

What does that return warrant?

What broader relation follows for the same reason?

What would falsify that generalization?

Under what guard does it remain valid?

What can now be excluded?

What can now be merged?

What must now be split?

What would make a folded distinction matter again?

Did this answer produce a nonredundant consequential change?

What consequential residual remains?

What are the two sides of that residual?

What relation distinguishes them?

What question now separates them?

Then recur.

\begin{center}\rule{0.5\linewidth}{0.5pt}\end{center}

\section{Part XXVIII --- Implementation
Layers}\label{part-xxviii-implementation-layers}

\subsection{106. Layer 1 --- Formal
controller}\label{layer-1-formal-controller}

Responsibilities:

\begin{itemize}
\tightlist
\item
  maintain inquiry state;
\item
  maintain obligations;
\item
  choose question contracts;
\item
  enforce type compatibility;
\item
  maintain claim graph;
\item
  maintain guard and warrant graph;
\item
  choose backends;
\item
  compute formal updates;
\item
  prevent candidate claims from becoming hard facts without promotion.
\end{itemize}

The formal controller should not depend on the language model
understanding the whole architecture.

\begin{center}\rule{0.5\linewidth}{0.5pt}\end{center}

\subsection{107. Layer 2 --- Question
compiler}\label{layer-2-question-compiler}

Input:

\[
O_t.
\]

Output:

\[
Q_\tau(args).
\]

Responsibilities:

\begin{itemize}
\tightlist
\item
  choose question role;
\item
  bind referents;
\item
  render plain-language form;
\item
  preserve role metadata;
\item
  resolve deixis such as ``this,'' ``that,'' and ``what changed?'' to
  explicit object IDs.
\end{itemize}

Surface wording may vary.

Formal question identity does not.

Two surface questions are equivalent when they instantiate the same
contract with the same bound arguments.

\begin{center}\rule{0.5\linewidth}{0.5pt}\end{center}

\subsection{108. Layer 3 --- Semantic
generator}\label{layer-3-semantic-generator}

The language model receives:

\begin{itemize}
\tightlist
\item
  minimal relevant local context;
\item
  one question or a dependency-safe batch;
\item
  referent material.
\end{itemize}

Its task is:

\[
\boxed{
\textbf{provide semantic content of the requested claim role.}
}
\]

It is not entrusted with:

\begin{itemize}
\tightlist
\item
  deciding what question comes next;
\item
  promoting its own answers to truth;
\item
  preserving formal consistency;
\item
  selecting its own warrant;
\item
  or changing protected consequences silently.
\end{itemize}

\begin{center}\rule{0.5\linewidth}{0.5pt}\end{center}

\subsection{109. Layer 4 --- Claim binder}\label{layer-4-claim-binder}

The binder creates:

\[
Claim_Q(a).
\]

It should perform minimal processing.

At the thinnest implementation, it needs only:

\begin{itemize}
\tightlist
\item
  question ID;
\item
  input object IDs;
\item
  raw answer text;
\item
  local scope;
\item
  provenance.
\end{itemize}

\begin{center}\rule{0.5\linewidth}{0.5pt}\end{center}

\subsection{110. Layer 5 --- Semantic
reifier}\label{layer-5-semantic-reifier}

Reification is invoked only when required.

Possible implementations:

\begin{itemize}
\tightlist
\item
  structured LLM extraction;
\item
  symbolic parser;
\item
  ontology mapper;
\item
  constraint compiler;
\item
  embedding-backed entity linker;
\item
  human confirmation;
\item
  domain adapter.
\end{itemize}

Failure is acceptable.

A failed reification creates another obligation rather than corrupting
state.

\begin{center}\rule{0.5\linewidth}{0.5pt}\end{center}

\subsection{111. Layer 6 --- Formal
backends}\label{layer-6-formal-backends}

Candidate adapters include:

\begin{itemize}
\tightlist
\item
  SAT solver;
\item
  SMT solver;
\item
  MaxSAT/OMT;
\item
  CHC solver;
\item
  PDR/IC3;
\item
  theorem prover;
\item
  symbolic executor;
\item
  planning engine;
\item
  simulator;
\item
  experiment runner;
\item
  database/retrieval system;
\item
  controller synthesizer;
\item
  external human or machine oracle.
\end{itemize}

Each backend exposes one common interface:

\[
\boxed{
\operatorname{Discharge}(O,e)\Downarrow r.
}
\]

\begin{center}\rule{0.5\linewidth}{0.5pt}\end{center}

\subsection{112. Layer 7 --- Warrant
checker}\label{layer-7-warrant-checker}

The warrant checker determines whether a backend return licenses
promotion.

Recommended warrant classes:

\begin{itemize}
\tightlist
\item
  formal proof;
\item
  checked certificate;
\item
  exhaustive computation;
\item
  authenticated observation;
\item
  experiment;
\item
  reproducible execution;
\item
  trusted retrieval;
\item
  corroborated evidence;
\item
  heuristic support;
\item
  model proposal.
\end{itemize}

Only classes authorized by policy may make hard cuts.

\begin{center}\rule{0.5\linewidth}{0.5pt}\end{center}

\section{Part XXIX --- Backend Adapter
Contract}\label{part-xxix-backend-adapter-contract}

\subsection{113. Generic adapter}\label{generic-adapter}

Every backend \(e\) implements:

\[
\boxed{
Encode_e
:
O
\to
Query_e
}
\]

\[
\boxed{
Run_e
:
Query_e
\to
Return_e
}
\]

\[
\boxed{
Decode_e
:
Return_e
\to
Result
}
\]

\[
\boxed{
Check_e
:
Result
\to
WarrantStatus.
}
\]

A result should expose:

\begin{Shaded}
\begin{Highlighting}[]
\NormalTok{Result \{}
\NormalTok{    kind}
\NormalTok{    proposition\_or\_relation}
\NormalTok{    scope}
\NormalTok{    assumptions}
\NormalTok{    witness\_or\_certificate}
\NormalTok{    provenance}
\NormalTok{    confidence\_or\_warrant\_class}
\NormalTok{\}}
\end{Highlighting}
\end{Shaded}

\begin{center}\rule{0.5\linewidth}{0.5pt}\end{center}

\section{Part XXX --- Return Types}\label{part-xxx-return-types}

\subsection{114. Canonical return
taxonomy}\label{canonical-return-taxonomy}

A backend return is normalized to one of:

\[
\boxed{
\begin{aligned}
&Witness\\
&Counterexample\\
&Separator\\
&EquivalenceCertificate\\
&Conflict\\
&Prerequisite\\
&ReachabilityWitness\\
&UnreachabilityCertificate\\
&Success\\
&Failure\\
&Unknown.
\end{aligned}
}
\]

The controller maps each return to an update and a next-obligation set.

\begin{center}\rule{0.5\linewidth}{0.5pt}\end{center}

\section{Part XXXI --- Conflict
Handling}\label{part-xxxi-conflict-handling}

\subsection{115. Candidate contradiction}\label{candidate-contradiction}

If candidate claims \(c_i,c_j\) conflict:

\[
\boxed{
Conflict(c_i,c_j)
}
\]

is added.

Do not:

\begin{itemize}
\tightlist
\item
  delete both automatically;
\item
  choose one arbitrarily;
\item
  infer arbitrary claims.
\end{itemize}

Generate:

\[
\boxed{
ResolveConflict(c_i,c_j).
}
\]

Possible resolutions:

\begin{itemize}
\tightlist
\item
  refute \(c_i\);
\item
  refute \(c_j\);
\item
  narrow one scope;
\item
  introduce hidden guard;
\item
  distinguish meanings;
\item
  discover invalid counterexample;
\item
  rebind representation;
\item
  retain unresolved conflict.
\end{itemize}

\begin{center}\rule{0.5\linewidth}{0.5pt}\end{center}

\subsection{116. Warrant conflict}\label{warrant-conflict}

If two active warranted lemmas produce inconsistency, treat this as a
higher-priority warrant failure.

Suspend any irreversible downstream step that depends on the
inconsistent conjunction until the conflict is localized.

Possible causes:

\begin{itemize}
\tightlist
\item
  invalid certificate;
\item
  incompatible scopes;
\item
  changed context;
\item
  unsound backend;
\item
  incorrect abstraction;
\item
  hidden assumption mismatch.
\end{itemize}

\begin{center}\rule{0.5\linewidth}{0.5pt}\end{center}

\section{Part XXXII --- Generalization}\label{part-xxxii-generalization}

\subsection{117. Generalization frontier}\label{generalization-frontier}

For return \(r\), let:

\[
\boxed{
Gen(r)
}
\]

be the nondominated frontier of warranted generalizations.

No unique strongest generalization is assumed.

Generalization may be implemented by:

\begin{itemize}
\tightlist
\item
  CDCL conflict analysis;
\item
  PDR inductive generalization;
\item
  Craig interpolation;
\item
  minimal support extraction;
\item
  MUS/MCS;
\item
  symbolic simplification;
\item
  theory-specific explanation;
\item
  contrast minimization.
\end{itemize}

Every generalized claim is itself attackable.

\begin{center}\rule{0.5\linewidth}{0.5pt}\end{center}

\section{Part XXXIII --- Compression and
Refinement}\label{part-xxxiii-compression-and-refinement}

\subsection{118. Refinement}\label{refinement}

If:

\[
\alpha(x)=\alpha(y)
\]

but:

\[
\Phi(x)\neq\Phi(y),
\]

the current abstraction is too coarse.

Seek:

\[
p(x)\neq p(y)
\]

and refine:

\[
B_{t+1}=B_t\cup\{p\}.
\]

\begin{center}\rule{0.5\linewidth}{0.5pt}\end{center}

\subsection{119. Compression}\label{compression}

A distinction \(p\in B_t\) is removable if deleting it preserves:

\begin{itemize}
\tightlist
\item
  protected consequence;
\item
  protected transition semantics;
\item
  protected control semantics.
\end{itemize}

The representation may therefore move both ways:

\[
\boxed{
\text{split when consequence requires distinction;}
}
\]

\[
\boxed{
\text{merge when consequence permits equivalence.}
}
\]

The representation is not monotonic.

The warrant store is.

\begin{center}\rule{0.5\linewidth}{0.5pt}\end{center}

\section{Part XXXIV --- Progress Law}\label{part-xxxiv-progress-law}

\subsection{120. Accepted progress}\label{accepted-progress}

A step counts as retained progress iff it produces at least one of:

\begin{enumerate}
\def\labelenumi{\arabic{enumi}.}
\tightlist
\item
  new nonredundant warranted conditional relation;
\item
  elimination of live possibility;
\item
  separator requiring refinement;
\item
  equivalence permitting compression;
\item
  localized consequence boundary;
\item
  falsified necessity;
\item
  falsified sufficiency;
\item
  established or eliminated prerequisite;
\item
  localized arrangement obstruction;
\item
  localized succession obstruction;
\item
  reopening condition;
\item
  changed consequential obligation.
\end{enumerate}

Otherwise the semantic output may be discarded.

Thus:

\[
\boxed{
\textbf{search may wander;
retained inquiry may not.}
}
\]

\begin{center}\rule{0.5\linewidth}{0.5pt}\end{center}

\section{Part XXXV --- Stopping Law}\label{part-xxxv-stopping-law}

\subsection{121. Stop when}\label{stop-when}

Stop an obligation when one of the following holds:

\subsubsection{Satisfied}\label{satisfied}

The desired claim is adequately warranted.

\subsubsection{Impossible}\label{impossible}

The target is adequately shown impossible within the declared binding.

\subsubsection{Consequence-equivalent
residual}\label{consequence-equivalent-residual}

All remaining alternatives yield identical protected consequence.

\subsubsection{No available
discriminator}\label{no-available-discriminator}

No representable question or available backend can distinguish the
remaining alternatives.

The final status is then:

\[
\boxed{
\text{unknown under the present binding}
}
\]

rather than an invented conclusion.

\begin{center}\rule{0.5\linewidth}{0.5pt}\end{center}

\section{Part XXXVI --- Minimum Viable
Implementation}\label{part-xxxvi-minimum-viable-implementation}

\subsection{122. Phase 1 --- Question-driven claim
graph}\label{phase-1-question-driven-claim-graph}

Implement first:

\begin{itemize}
\tightlist
\item
  typed question contracts;
\item
  claim wrapper;
\item
  bound referents;
\item
  candidate graph;
\item
  conflict edges;
\item
  recursive question selection;
\item
  no hard formalization requirement.
\end{itemize}

Goal:

prove that lawful question composition can drive an LLM through coherent
recursive inquiry while tolerating bad answers.

\begin{center}\rule{0.5\linewidth}{0.5pt}\end{center}

\subsection{123. Phase 2 --- Lazy
formalization}\label{phase-2-lazy-formalization}

Add:

\begin{itemize}
\tightlist
\item
  L0--L3 semantic levels;
\item
  reification adapters;
\item
  SAT/SMT predicates where possible;
\item
  explicit necessity/sufficiency checks;
\item
  scope and guard extraction.
\end{itemize}

\begin{center}\rule{0.5\linewidth}{0.5pt}\end{center}

\subsection{124. Phase 3 --- Abstract
state}\label{phase-3-abstract-state}

Add:

\begin{itemize}
\tightlist
\item
  predicate basis;
\item
  abstraction map;
\item
  separator-driven refinement;
\item
  consequence-relative compression;
\item
  CEGAR-style spurious-counterexample handling.
\end{itemize}

\begin{center}\rule{0.5\linewidth}{0.5pt}\end{center}

\subsection{125. Phase 4 --- Recursive formal
engine}\label{phase-4-recursive-formal-engine}

Add:

\begin{itemize}
\tightlist
\item
  CHC obligation compilation;
\item
  Spacer/PDR;
\item
  predecessor propagation;
\item
  inductive cuts;
\item
  generalized failure learning.
\end{itemize}

\begin{center}\rule{0.5\linewidth}{0.5pt}\end{center}

\subsection{126. Phase 5 --- Control}\label{phase-5-control}

Add:

\begin{itemize}
\tightlist
\item
  action relations;
\item
  may/must predecessors;
\item
  controller synthesis;
\item
  feedback refinement;
\item
  concrete execution verification.
\end{itemize}

\begin{center}\rule{0.5\linewidth}{0.5pt}\end{center}

\subsection{127. Phase 6 --- Multi-backend
warrant}\label{phase-6-multi-backend-warrant}

Add:

\begin{itemize}
\tightlist
\item
  proof-producing SMT;
\item
  retrieval adapters;
\item
  observational evidence;
\item
  experiment adapters;
\item
  provenance and certificate checker.
\end{itemize}

\begin{center}\rule{0.5\linewidth}{0.5pt}\end{center}

\section{Part XXXVII --- Data
Structures}\label{part-xxxvii-data-structures}

\subsection{128. QuestionContract}\label{questioncontract}

\begin{Shaded}
\begin{Highlighting}[]
\NormalTok{QuestionContract \{}
\NormalTok{    id}
\NormalTok{    input\_roles[]}
\NormalTok{    output\_claim\_role}
\NormalTok{    precondition}
\NormalTok{    render\_template}
\NormalTok{    bind\_policy}
\NormalTok{    reifier}
\NormalTok{    verifier}
\NormalTok{    update\_rule}
\NormalTok{    next\_obligation\_rules[]}
\NormalTok{\}}
\end{Highlighting}
\end{Shaded}

\begin{center}\rule{0.5\linewidth}{0.5pt}\end{center}

\subsection{129. Claim}\label{claim}

\begin{Shaded}
\begin{Highlighting}[]
\NormalTok{Claim \{}
\NormalTok{    id}
\NormalTok{    role}
\NormalTok{    bound\_args[]}
\NormalTok{    payload}
\NormalTok{    scope}
\NormalTok{    provenance}
\NormalTok{    status}
\NormalTok{    representation\_level}
\NormalTok{    parents[]}
\NormalTok{    conflicts[]}
\NormalTok{\}}
\end{Highlighting}
\end{Shaded}

\begin{center}\rule{0.5\linewidth}{0.5pt}\end{center}

\subsection{130. FormalCandidate}\label{formalcandidate}

\begin{Shaded}
\begin{Highlighting}[]
\NormalTok{FormalCandidate \{}
\NormalTok{    claim\_id}
\NormalTok{    formal\_type}
\NormalTok{    expression}
\NormalTok{    carrier}
\NormalTok{    scope}
\NormalTok{    assumptions[]}
\NormalTok{\}}
\end{Highlighting}
\end{Shaded}

\begin{center}\rule{0.5\linewidth}{0.5pt}\end{center}

\subsection{131. WarrantedLemma}\label{warrantedlemma}

\begin{Shaded}
\begin{Highlighting}[]
\NormalTok{WarrantedLemma \{}
\NormalTok{    id}
\NormalTok{    guard}
\NormalTok{    relation}
\NormalTok{    scope}
\NormalTok{    warrant\_class}
\NormalTok{    dependencies[]}
\NormalTok{    certificate}
\NormalTok{    source\_claims[]}
\NormalTok{\}}
\end{Highlighting}
\end{Shaded}

\begin{center}\rule{0.5\linewidth}{0.5pt}\end{center}

\subsection{132. Obligation}\label{obligation}

\begin{Shaded}
\begin{Highlighting}[]
\NormalTok{Obligation \{}
\NormalTok{    id}
\NormalTok{    kind}
\NormalTok{    carrier}
\NormalTok{    args[]}
\NormalTok{    scope}
\NormalTok{    parent\_obligations[]}
\NormalTok{    priority\_vector}
\NormalTok{    status}
\NormalTok{\}}
\end{Highlighting}
\end{Shaded}

\begin{center}\rule{0.5\linewidth}{0.5pt}\end{center}

\subsection{133. InquiryState}\label{inquirystate}

\begin{Shaded}
\begin{Highlighting}[]
\NormalTok{InquiryState \{}
\NormalTok{    binding}
\NormalTok{    context}
\NormalTok{    protected\_consequences}
\NormalTok{    candidate\_claim\_graph}
\NormalTok{    warranted\_lemma\_store}
\NormalTok{    predicate\_basis}
\NormalTok{    abstraction}
\NormalTok{    abstract\_transitions}
\NormalTok{    obligations}
\NormalTok{    warrant\_graph}
\NormalTok{    current\_mode}
\NormalTok{\}}
\end{Highlighting}
\end{Shaded}

\begin{center}\rule{0.5\linewidth}{0.5pt}\end{center}

\section{Part XXXVIII --- System
Invariants}\label{part-xxxviii-system-invariants}

\subsection{134. Invariant I --- Claim/fact
separation}\label{invariant-i-claimfact-separation}

\[
\boxed{
Claim[\tau]\not\Rightarrow Warranted[\tau].
}
\]

\begin{center}\rule{0.5\linewidth}{0.5pt}\end{center}

\subsection{135. Invariant II --- No candidate
explosion}\label{invariant-ii-no-candidate-explosion}

Contradictory candidate claims create conflicts, not arbitrary
entailments.

\begin{center}\rule{0.5\linewidth}{0.5pt}\end{center}

\subsection{136. Invariant III --- Guarded
retention}\label{invariant-iii-guarded-retention}

Every retained hard relation has:

\begin{itemize}
\tightlist
\item
  guard;
\item
  scope;
\item
  dependencies;
\item
  warrant.
\end{itemize}

\begin{center}\rule{0.5\linewidth}{0.5pt}\end{center}

\subsection{137. Invariant IV --- No hidden control
promotion}\label{invariant-iv-no-hidden-control-promotion}

Descriptive factorization cannot itself create a control certificate.

\begin{center}\rule{0.5\linewidth}{0.5pt}\end{center}

\subsection{138. Invariant V --- Counterexample
pressure}\label{invariant-v-counterexample-pressure}

Every claimed necessity, sufficiency, and prerequisite creates the
corresponding counterexample obligation unless already discharged.

\begin{center}\rule{0.5\linewidth}{0.5pt}\end{center}

\subsection{139. Invariant VI --- Consequence-relative
representation}\label{invariant-vi-consequence-relative-representation}

No active distinction may be compressed if a protected consequence,
transition, or control relation can inspect it.

\begin{center}\rule{0.5\linewidth}{0.5pt}\end{center}

\subsection{140. Invariant VII --- Warrant
DAG}\label{invariant-vii-warrant-dag}

Positive warrant remains acyclic.

\begin{center}\rule{0.5\linewidth}{0.5pt}\end{center}

\subsection{141. Invariant VIII --- Unknown is
legal}\label{invariant-viii-unknown-is-legal}

Failure to distinguish is not converted into equivalence without
warrant.

Failure to find a counterexample is not converted into necessity.

Failure to actualize is not converted into impossibility without
adequate warrant.

\begin{center}\rule{0.5\linewidth}{0.5pt}\end{center}

\section{Part XXXIX --- Verification
Tests}\label{part-xxxix-verification-tests}

\subsection{142. Protocol tests}\label{protocol-tests}

A conforming implementation should pass the following.

\subsubsection{Arbitrary-payload safety}\label{arbitrary-payload-safety}

Feed arbitrary strings as answers to every question contract.

Expected result:

\begin{itemize}
\tightlist
\item
  candidate claims are created;
\item
  no crash from semantic ill-typing;
\item
  no automatic promotion to hard theory.
\end{itemize}

\subsubsection{Contradiction
containment}\label{contradiction-containment}

Create claims for \(P\) and \(\neg P\).

Expected result:

\begin{itemize}
\tightlist
\item
  conflict obligation;
\item
  no arbitrary derived claim.
\end{itemize}

\subsubsection{Necessity attack
completeness}\label{necessity-attack-completeness}

Whenever a claim is tagged as necessity, verify that an obligation
equivalent to ``consequence without necessity'' exists or is already
discharged.

\subsubsection{Sufficiency attack
completeness}\label{sufficiency-attack-completeness}

Whenever a claim is tagged as sufficient, verify that its opposite cell
is queried or discharged.

\subsubsection{Guard reopening}\label{guard-reopening}

Activate a lemma, invalidate its guard, and verify:

\begin{itemize}
\tightlist
\item
  lemma remains stored;
\item
  it ceases to constrain the active model;
\item
  dependent distinctions may reopen.
\end{itemize}

\subsubsection{Conservative refinement}\label{conservative-refinement}

Add a predicate and verify the old abstraction is recoverable through
forgetting.

\subsubsection{Compression safety}\label{compression-safety}

Attempt to merge consequence-distinct states and verify refusal.

\subsubsection{Description/control
separation}\label{descriptioncontrol-separation}

Supply a perfect predictor with no implementable transition and verify
that no control certificate is produced.

\begin{center}\rule{0.5\linewidth}{0.5pt}\end{center}

\section{Part XL --- Experimental Program for LLM
Use}\label{part-xl-experimental-program-for-llm-use}

\subsection{143. Delivery regimes}\label{delivery-regimes}

Evaluate at least:

\begin{enumerate}
\def\labelenumi{\arabic{enumi}.}
\tightlist
\item
  full theory in one prompt;
\item
  all questions in one checklist;
\item
  dependency-aware groups;
\item
  one adaptive question at a time;
\item
  rhythmic adaptive questioning.
\end{enumerate}

Measure:

\begin{itemize}
\tightlist
\item
  contradiction rate;
\item
  false necessity rate;
\item
  correction rate;
\item
  final consequence accuracy;
\item
  number of retained claims;
\item
  average claim lifetime;
\item
  amount of context compression;
\item
  question-order sensitivity;
\item
  counterexample discovery;
\item
  semantic drift;
\item
  warrant quality.
\end{itemize}

\begin{center}\rule{0.5\linewidth}{0.5pt}\end{center}

\subsection{144. Expected hypothesis}\label{expected-hypothesis}

The strongest predicted regime is:

\[
\boxed{
\text{one locally typed question at a time,
chosen from the actual previous return,
with external state retention and compression}.
}
\]

This minimizes instruction averaging and forces downstream questions to
depend on actually produced claims rather than imagined intermediate
conclusions.

\begin{center}\rule{0.5\linewidth}{0.5pt}\end{center}

\section{Part XLI --- Research
Questions}\label{part-xli-research-questions}

\subsection{145. Open technical
questions}\label{open-technical-questions}

\subsubsection{Question scheduling}\label{question-scheduling}

What objective best predicts the long-run value of a question?

Possibilities include:

\begin{itemize}
\tightlist
\item
  version-space reduction;
\item
  expected generalized-cut size;
\item
  expected boundary localization;
\item
  proof-obligation depth;
\item
  risk-adjusted information gain.
\end{itemize}

\subsubsection{Semantic reification}\label{semantic-reification}

How little parsing is sufficient before formal machinery becomes useful?

\subsubsection{Claim-role robustness}\label{claim-role-robustness}

How sensitive is inquiry quality to surface wording when formal question
role remains fixed?

\subsubsection{Context compression}\label{context-compression}

How aggressively can candidate history be compressed while preserving
future reopening conditions?

\subsubsection{Contradiction
localization}\label{contradiction-localization}

Can conflict graphs over typed opaque claims reliably identify which
semantic claims should be re-presented to the LLM for repair?

\subsubsection{Order commutativity}\label{order-commutativity}

Can question permutation distinguish genuine relational interaction from
transformer anchoring effects?

\subsubsection{Cross-domain transfer}\label{cross-domain-transfer}

Does the same question calculus induce useful inquiry in domains with
radically different native formalisms?

\subsubsection{Formal completeness}\label{formal-completeness}

Under what restrictions on \(\mathcal L,\mathcal E,\mathcal U\) does RCI
become complete for a given class of inquiry problems?

\begin{center}\rule{0.5\linewidth}{0.5pt}\end{center}

\section{Part XLII --- Governing Laws}\label{part-xlii-governing-laws}

\subsection{146. Generation law}\label{generation-law}

\[
\boxed{
\textbf{Generate for maximal consequential discrimination,
not for plausibility.}
}
\]

\begin{center}\rule{0.5\linewidth}{0.5pt}\end{center}

\subsection{147. Boundary law}\label{boundary-law}

\[
\boxed{
\textbf{Within a consequence class, vary aggressively.
Across consequence classes, change minimally.}
}
\]

\begin{center}\rule{0.5\linewidth}{0.5pt}\end{center}

\subsection{148. Falsification law}\label{falsification-law}

\[
\boxed{
\textbf{Every “must” generates “can the result occur without it?”}
}
\]

\[
\boxed{
\textbf{Every “enough” generates “can it occur while the result fails?”}
}
\]

\begin{center}\rule{0.5\linewidth}{0.5pt}\end{center}

\subsection{149. Prerequisite law}\label{prerequisite-law}

\[
\boxed{
\textbf{Every proposed prerequisite is attacked by attempting the target without it.}
}
\]

\begin{center}\rule{0.5\linewidth}{0.5pt}\end{center}

\subsection{150. Control law}\label{control-law}

\[
\boxed{
\textbf{Never infer controllability from descriptive correlation or factorization.}
}
\]

\begin{center}\rule{0.5\linewidth}{0.5pt}\end{center}

\subsection{151. Question law}\label{question-law}

\[
\boxed{
\textbf{Questions create typed claims, not facts.}
}
\]

\begin{center}\rule{0.5\linewidth}{0.5pt}\end{center}

\subsection{152. Fallibility law}\label{fallibility-law}

\[
\boxed{
\textbf{The question graph must remain lawful even when semantic answers are wrong.}
}
\]

\begin{center}\rule{0.5\linewidth}{0.5pt}\end{center}

\subsection{153. Conflict law}\label{conflict-law}

\[
\boxed{
\textbf{Contradiction creates an obligation; it does not create explosion.}
}
\]

\begin{center}\rule{0.5\linewidth}{0.5pt}\end{center}

\subsection{154. Warrant law}\label{warrant-law}

\[
\boxed{
\textbf{Only independently licensed warrant may promote a provisional claim into retained hard knowledge.}
}
\]

\begin{center}\rule{0.5\linewidth}{0.5pt}\end{center}

\subsection{155. Compression law}\label{compression-law}

\[
\boxed{
\textbf{Represent no distinction more finely than protected consequence requires,
and no more coarsely than protected consequence permits.}
}
\]

\begin{center}\rule{0.5\linewidth}{0.5pt}\end{center}

\subsection{156. Reopening law}\label{reopening-law}

\[
\boxed{
\textbf{A changed guard deactivates a conclusion; it need not erase the warrant that established it under its original conditions.}
}
\]

\begin{center}\rule{0.5\linewidth}{0.5pt}\end{center}

\subsection{157. Progress law}\label{progress-law}

\[
\boxed{
\textbf{Accepted reasoning must create a nonredundant consequential change.}
}
\]

\begin{center}\rule{0.5\linewidth}{0.5pt}\end{center}

\subsection{158. Recursion law}\label{recursion-law}

\[
\boxed{
\textbf{Every consequential residual becomes the next distinction.}
}
\]

\begin{center}\rule{0.5\linewidth}{0.5pt}\end{center}

\section{Part XLIII --- Complete Compressed
Definition}\label{part-xliii-complete-compressed-definition}

Ratcheting Consequence Inquiry is a typed recursive inquiry process over
relational state.

It begins from a protected consequence distinction rather than requiring
a pre-specified concrete solution.

It generates strongly differing realizations within consequence classes
to destroy accidental commonality and minimally differing realizations
across consequence classes to expose consequence boundaries.

It proposes compact noncircular relations through which protected
consequences may factor, then actively attempts to falsify every
proposed necessity and sufficiency by constructing the corresponding
missing counterexample cells.

It distinguishes descriptive consequence structure from actionable
control structure and recursively asks what must be established,
removed, repaired, or transformed for the desired consequence relation
to be actualized from the present arrangement.

Every proposed prerequisite is itself subjected to alternate-route
falsification.

The semantic generator is controlled through typed natural-language
question contracts.

A question determines the relation being requested and creates a
provisional claim of that type. The semantic answer need not initially
be formally parsed or trusted.

Candidate claims are allowed to conflict. Contradictions remain
localized and generate further inquiry rather than causing logical
explosion.

Claims are formalized only as needed. Independent proof, computation,
observation, experiment, execution, retrieval, or other licensed
discharge mechanisms may promote a candidate into a guarded warranted
relation.

Warranted relations accumulate monotonically in a conditional theory,
while active possibilities and representations may split, merge, reopen,
or deactivate as contexts and guards change.

The representation continuously refines whenever protected consequences
distinguish states currently identified and compresses whenever
protected consequences no longer require a distinction.

Every successful or failed inquiry return is localized and generalized
into the broadest warranted reusable relation available.

Every consequential residual recursively becomes the next inquiry.

In its minimal operational form:

\[
\boxed{
\textbf{
STRETCH WITHIN
\;\leftrightarrow\;
SQUEEZE ACROSS
\;\to\;
FACTOR
\;\to\;
FALSIFY
\;\to\;
ACTUALIZE
\;\to\;
DISCHARGE
\;\to\;
GENERALIZE
\;\to\;
LEARN CONDITIONALLY
\;\to\;
REFOLD
\;\to\;
RECUR.
}
}
\]

In its minimal language form:

\[
\boxed{
\begin{aligned}
&\textbf{What difference matters?}\\
&\textbf{What can change without changing it?}\\
&\textbf{What smallest change would reverse it?}\\
&\textbf{Can the result hold without what appears necessary?}\\
&\textbf{Can what appears sufficient hold while the result fails?}\\
&\textbf{What smaller relation survives those attacks?}\\
&\textbf{What would have to be true for that relation to hold here?}\\
&\textbf{Can the target be reached without each proposed prerequisite?}\\
&\textbf{What independently constrains the answer?}\\
&\textbf{What does the return actually warrant?}\\
&\textbf{What distinction can now be forgotten?}\\
&\textbf{What would make it matter again?}\\
&\textbf{What consequential difference remains?}
\end{aligned}
}
\]

The answer to the final question becomes the referent of the first
question in the next recursion.

\begin{center}\rule{0.5\linewidth}{0.5pt}\end{center}

\section{References}\label{references}

{[}R1{]} Dexter Kozen. \emph{Kleene Algebra with Tests}. ACM
Transactions on Programming Languages and Systems, 19(3):427--443, 1997.
DOI: 10.1145/256167.256195.

{[}R2{]} Dana Angluin. \emph{Learning Regular Sets from Queries and
Counterexamples}. Information and Computation, 75:87--106, 1987.

{[}R3{]} Robert Paige and Robert E. Tarjan. \emph{Three Partition
Refinement Algorithms}. SIAM Journal on Computing, 16(6):973--989, 1987.
DOI: 10.1137/0216062.

{[}R4{]} Patrick Cousot and Radhia Cousot. \emph{Abstract
Interpretation: A Unified Lattice Model for Static Analysis of Programs
by Construction or Approximation of Fixpoints}. POPL, pp.~238--252,
1977.

{[}R5{]} Robert Nieuwenhuis, Albert Oliveras, and Cesare Tinelli.
\emph{Solving SAT and SAT Modulo Theories: From an Abstract DPLL
Procedure to DPLL(T).} Journal of the ACM, 53(6):937--977, 2006. DOI:
10.1145/1217856.1217859.

{[}R6{]} Gunther Reissig, Alexander Weber, and Matthias Rungger.
\emph{Feedback Refinement Relations for the Synthesis of Symbolic
Controllers}. 2015/2017.

{[}R7{]} Ulrich Junker. \emph{QUICKXPLAIN: Preferred Explanations and
Relaxations for Over-Constrained Problems}. AAAI, pp.~167--172, 2004.

{[}R8{]} Johan de Kleer. \emph{An Assumption-Based TMS}. Artificial
Intelligence, 28(2):127--162, 1986. DOI: 10.1016/0004-3702(86)90080-9.

{[}R9{]} Edmund M. Clarke, Orna Grumberg, Somesh Jha, Yuan Lu, and
Helmut Veith. \emph{Counterexample-Guided Abstraction Refinement}. CAV,
2000.

{[}R10{]} LFSC / SMT proof-checking work: logical frameworks with side
conditions support small trusted proof-checking kernels for SMT
certificates.

{[}R11{]} Aaron R. Bradley. \emph{SAT-Based Model Checking without
Unrolling}. VMCAI, pp.~70--87, 2011.

{[}R12{]} Z3 Spacer / generalized Property Directed Reachability for
constrained Horn clauses.

{[}R13{]} Kenneth L. McMillan. \emph{Interpolation and SAT-Based Model
Checking}. CAV, pp.~1--13, 2003. DOI: 10.1007/978-3-540-45069-6\_1.

\begin{center}\rule{0.5\linewidth}{0.5pt}\end{center}

\section{Project Definition}\label{project-definition}

\[
\boxed{
\textbf{RATCHETING CONSEQUENCE INQUIRY}
}
\]

is therefore defined as:

A proof-producing, property-directed, question-driven refinement process
over an adaptive relational abstraction in which natural-language
questions create typed provisional claims, semantic generators fill
those claims with content, lawful composition exposes boundaries and
contradictions, independent discharge mechanisms warrant reusable
relations, and every consequential residual recursively generates the
next inquiry.

Its architectural principle is:

\[
\boxed{
\textbf{
FORMAL OBLIGATION
\to
PLAIN-LANGUAGE QUESTION
\to
UNTRUSTED SEMANTIC CLAIM
\to
LAWFUL COMPOSITION
\to
FORMAL/EXTERNAL DISCHARGE
\to
GUARDED WARRANT
\to
REFINEMENT/COMPRESSION
\to
NEXT OBLIGATION.
}
}
\]

Its implementation criterion is:

\[
\boxed{
\textbf{The process must remain formally well-defined even when the semantic generator is wrong.}
}
\]

Its epistemic criterion is:

\[
\boxed{
\textbf{Only independently warranted consequential changes enter retained knowledge.}
}
\]

Its recursive criterion is:

\[
\boxed{
\textbf{Every unresolved consequential difference becomes another question.}
}
\]

\end{document}
